%=============================================================
% LIBRO: EL SECRETO DE LOS RETICULOS
% Criptografia Post-Cuantica explicada para estudiantes
% de primer semestre de matematicas basica
%=============================================================
\documentclass[11pt,a4paper,oneside]{book}

% Lenguaje y codificacion
\usepackage[utf8]{inputenc}
\usepackage[T1]{fontenc}
\usepackage[spanish,es-noquoting,es-tabla]{babel}
\usepackage{lmodern}
\usepackage{microtype}

% Matematicas
\usepackage{amsmath,amssymb,amsthm,mathtools}
\usepackage{bm}

% Geometria de pagina
\usepackage[
  a4paper,
  top=2.8cm, bottom=2.8cm,
  left=2.8cm, right=2.8cm,
  headheight=16pt
]{geometry}

% Encabezados elegantes
\usepackage{fancyhdr}
\pagestyle{fancy}
\fancyhf{}
\fancyhead[L]{\small\sffamily\color{gray!80!black}\nouppercase{\leftmark}}
\fancyfoot[C]{\small\sffamily\thepage}
\renewcommand{\headrulewidth}{0.5pt}
\renewcommand{\footrulewidth}{0pt}

% Colores
\usepackage[dvipsnames,svgnames]{xcolor}
\definecolor{PQCMainBlue}{RGB}{24, 76, 120}
\definecolor{PQCDark}{RGB}{15, 30, 50}
\definecolor{PQCTeal}{RGB}{12, 133, 143}
\definecolor{PQCGreen}{RGB}{34, 139, 84}
\definecolor{PQCAmber}{RGB}{217, 119, 6}
\definecolor{PQCPurple}{RGB}{109, 40, 217}
\definecolor{PQCCrimson}{RGB}{185, 28, 28}
\definecolor{BoxBgBlue}{RGB}{243, 248, 253}
\definecolor{BoxBgGreen}{RGB}{240, 253, 244}
\definecolor{BoxBgAmber}{RGB}{254, 252, 232}
\definecolor{BoxBgPurple}{RGB}{248, 245, 255}
\definecolor{BoxBgTeal}{RGB}{240, 253, 250}

% TikZ
\usepackage{tikz}
\usetikzlibrary{arrows.meta, positioning, calc, shapes.geometric, patterns, decorations.pathreplacing, backgrounds}

% Cajas didacticas (tcolorbox)
\usepackage{tcolorbox}
\tcbuselibrary{breakable, skins, theorems}

% 1. Caja de Intuicion (Purpura)
\newtcolorbox{intuicion}[1][]{
  colback=BoxBgPurple, colframe=PQCPurple,
  fonttitle=\bfseries\sffamily\color{white},
  coltitle=white, colbacktitle=PQCPurple,
  title={$\blacktriangleright$ Intuici\'{o}n Clave: #1},
  breakable, enhanced, arc=2.5mm, boxrule=0.9pt,
  drop fuzzy shadow=black!12, top=3mm, bottom=3mm, left=3mm, right=3mm
}

% 2. Caja de Metafora Cotidiana (Verde)
\newtcolorbox{metafora}[1][]{
  colback=BoxBgGreen, colframe=PQCGreen,
  fonttitle=\bfseries\sffamily\color{white},
  coltitle=white, colbacktitle=PQCGreen,
  title={$\bigstar$ Met\'{a}fora Cotidiana: #1},
  breakable, enhanced, arc=2.5mm, boxrule=0.9pt,
  drop fuzzy shadow=black!12, top=3mm, bottom=3mm, left=3mm, right=3mm
}

% 3. Caja de Ejemplo Paso a Paso (Azul / Teal)
\newtcolorbox{ejemplonumerico}[1][]{
  colback=BoxBgTeal, colframe=PQCTeal,
  fonttitle=\bfseries\sffamily\color{white},
  coltitle=white, colbacktitle=PQCTeal,
  title={$\mathbb{Z}_q$ Paso a Paso: #1},
  breakable, enhanced, arc=2.5mm, boxrule=0.9pt,
  drop fuzzy shadow=black!12, top=3mm, bottom=3mm, left=3mm, right=3mm
}

% 4. Caja de Ojo / Alerta didactica (Ambar)
\newtcolorbox{ojo}[1][]{
  colback=BoxBgAmber, colframe=PQCAmber,
  fonttitle=\bfseries\sffamily\color{black!90},
  coltitle=black!90, colbacktitle=PQCAmber!60!white,
  title={$\triangle$ ¡Atenci\'{o}n con esto!: #1},
  breakable, enhanced, arc=2.5mm, boxrule=0.9pt,
  drop fuzzy shadow=black!12, top=3mm, bottom=3mm, left=3mm, right=3mm
}

% 5. Caja de Resumen del Capitulo (Azul Marino)
\newtcolorbox{resumencapitulo}[1][Ideas que debes llevarte en tu mochila]{
  colback=BoxBgBlue, colframe=PQCMainBlue,
  fonttitle=\bfseries\sffamily\color{white},
  coltitle=white, colbacktitle=PQCMainBlue,
  title={\checkmark\ #1},
  breakable, enhanced, arc=2.5mm, boxrule=1pt,
  drop fuzzy shadow=black!15, top=3.5mm, bottom=3.5mm, left=4mm, right=4mm
}

% Tablas
\usepackage{booktabs}
\usepackage{array}

% Enlaces
\usepackage{hyperref}
\hypersetup{
  colorlinks=true,
  linkcolor=PQCMainBlue,
  citecolor=PQCTeal,
  urlcolor=PQCPurple,
  pdftitle={El Secreto de los Reticulos - Criptografia Post-Cuantica para Primer Semestre},
  pdfauthor={Universidad Mayor de San Andres - UMSA}
}

% Formato de titulos de capitulo
\usepackage{titlesec}
\titleformat{\chapter}[display]
  {\normalfont\sffamily\huge\bfseries\color{PQCMainBlue}}
  {\filleft\small\scshape\color{gray}Cap\'itulo \thechapter}
  {0.5ex}
  {\vspace{0.5ex}\hrule height 1.2pt\vspace{1.5ex}}
  [\vspace{1.5ex}\hrule height 0.5pt]

\titleformat{\section}
  {\normalfont\sffamily\Large\bfseries\color{PQCDark}}
  {\thesection}{1em}{}

\titleformat{\subsection}
  {\normalfont\sffamily\large\bfseries\color{PQCTeal}}
  {\thesubsection}{1em}{}

% Comandos utiles
\newcommand{\Z}{\mathbb{Z}}
\newcommand{\R}{\mathbb{R}}
\newcommand{\Zq}{\mathbb{Z}_q}
\newcommand{\Rq}{R_q}

%=============================================================
\begin{document}

% Portada
\begin{titlepage}
\begin{tikzpicture}[remember picture, overlay]
  % Fondo decorativo
  \fill[BoxBgBlue] (current page.south west) rectangle (current page.north east);
  
  % Franja lateral de color
  \fill[PQCMainBlue] (current page.north west) rectangle ($(current page.south west)+(2.0cm,0)$);
  \fill[PQCTeal] ($(current page.north west)+(2.0cm,0)$) rectangle ($(current page.south west)+(2.2cm,0)$);
  
  % Reticulo sutil en el fondo de la portada
  \begin{scope}[shift={($(current page.north east)+(-3,-5)$)}, opacity=0.15]
    \foreach \x in {-5,...,4} {
      \foreach \y in {-5,...,4} {
        \fill[PQCMainBlue] ({\x*1.2 + \y*0.4}, {\y*1.0}) circle (2pt);
      }
    }
  \end{scope}
\end{tikzpicture}

\vspace*{1.5cm}
\begin{flushleft}
  \hspace*{1.8cm}
  \begin{minipage}{0.88\textwidth}
    {\large\sffamily\bfseries\color{PQCTeal} UNIVERSIDAD MAYOR DE SAN ANDRÉS}\\[0.2cm]
    {\small\sffamily\color{gray} Facultad de Ciencias Puras y Naturales --- Carrera de Informática}\\[2.5cm]
    
    {\fontsize{30}{36}\sffamily\bfseries\color{PQCMainBlue} EL SECRETO DE LOS RETÍCULOS}\\[0.6cm]
    {\LARGE\sffamily\color{PQCPurple} Una Aventura Visual e Intuitiva por la Criptografía Post-Cuántica}\\[0.4cm]
    {\large\sffamily\color{gray!90!black} Conceptos, ideas y geometría moderna explicados para estudiantes de primer semestre}\\[2.5cm]
    
    \begin{tcolorbox}[colback=white, colframe=PQCTeal, arc=2mm, width=0.95\textwidth]
      \sffamily\small
      \textbf{¿Qué aprenderás aquí?} \\
      Aprenderás cómo los matemáticos protegen hoy la información del mundo utilizando cuadrículas geométricas multidimensionales (retículos), la aritmética de los relojes y un toque mágico de ``ruido'' matemático, sin enredos técnicos incomprensibles.
    \end{tcolorbox}
    
    \vfill
    \vspace*{2cm}
    {\small\sffamily\color{gray}
    \textbf{Edición Didáctica} --- La Paz, Bolivia --- 2026}
  \end{minipage}
\end{flushleft}
\end{titlepage}

% Pagina de creditos
\thispagestyle{empty}
\vspace*{\fill}
\begin{center}
\sffamily\small
\textbf{El Secreto de los Retículos}\\
\textit{Criptografía Post-Cuántica explicada para estudiantes de primer semestre}\\
\vspace{0.4cm}
Proyecto de Grado e Innovación Pedagógica\\
Facultad de Ciencias Puras y Naturales\\
Universidad Mayor de San Andrés (UMSA)\\
La Paz, Bolivia, 2026\\
\vspace{0.8cm}
\begin{minipage}{0.8\textwidth}
\centering\itshape\color{gray!80!black}
``La matemática no se trata de hacer cosas difíciles de manera complicada, sino de encontrar ideas simples que expliquen cosas maravillosas.''
\end{minipage}
\end{center}
\vspace*{\fill}
\newpage

% Tabla de contenidos
\tableofcontents
\newpage

% Capitulos
\chapter*{Prólogo: Para el Estudiante Curioso}
\addcontentsline{toc}{chapter}{Prólogo: Para el Estudiante Curioso}
\markboth{Prólogo}{Prólogo}

¡Te damos la bienvenida al emocionante mundo de las matemáticas universitarias!

Si estás en tu primer semestre de la carrera de Matemáticas, Informática o Ingeniería, es muy probable que tus clases giren en torno a conceptos que parecen sencillos a primera vista: vectores en el plano cartesiano $(x, y)$, sistemas de ecuaciones con dos incógnitas, productos de polinomios y divisiones con residuo.

A menudo, los libros de texto tradicionales presentan estos temas como una serie de recetas mecánicas: ``iguala a cero'', ``despeja la $x$'', ``multiplica fila por columna''. Sin embargo, detrás de esas operaciones elementales se esconde una de las revoluciones científicas y tecnológicas más fascinantes de nuestro siglo: \textbf{la Criptografía Post-Cuántica (PQC)}.

\begin{metafora}[El candado del futuro]
Imagina que durante décadas el mundo entero cerró sus puertas y secretos bancarios con candados basados en combinaciones numéricas que se pueden abrir si encuentras un ritmo o un eco repetitivo. Si alguien construye una máquina capaz de escuchar todos los ecos a la vez (una computadora cuántica), todos esos viejos candados se abrirán al instante.

¿Qué hicieron los matemáticos? No intentaron hacer las cerraduras viejas más grandes. En su lugar, inventaron un tipo de candado completamente nuevo: un laberinto geométrico infinito en cientos de dimensiones donde no existen ecos ni ritmos repetitivos. Para entrar, necesitas conocer un atajo secreto que solo tú posees.
\end{metafora}

\section*{¿De qué trata este libro?}

Este libro fue escrito pensando exclusivamente en ti. No asume que conoces teoría de anillos avanzada, ni análisis funcional, ni geometría diferencial. Todo lo que necesitas para entender las ideas maestras de la nueva era de la seguridad informática cabe en las herramientas de tu primer semestre:

\begin{enumerate}
  \item \textbf{Geometría intuitiva:} Saber qué es un punto en una cuadrícula y cómo sumar dos flechas (vectores).
  \item \textbf{Aritmética del reloj:} Comprender cómo funcionan las horas en un reloj de pared (donde después de las 12 viene el 1 otra vez).
  \item \textbf{Un toque de álgebra elemental:} Saber qué es un polinomio simple y cómo sumar $2x + 3x = 5x$.
\end{enumerate}

\section*{¿Cómo está organizado el viaje?}

El recorrido está dividido en siete etapas claras y progresivas:

\begin{itemize}
  \item \textbf{Capítulo 1:} Nos preguntaremos cómo se guarda un secreto en el universo de las matemáticas y por qué una simple ecuación lineal deja de ser fácil si le agregamos un ``ruido'' milimétrico.
  \item \textbf{Capítulo 2:} Exploraremos los \textbf{retículos} (\textit{lattices}), una generalización fascinante de la hoja cuadriculada, y descubriremos el dilema entre una base ``buena'' y una base ``mala''.
  \item \textbf{Capítulo 3:} Aprenderemos el problema central que domina la criptografía moderna: \textbf{LWE} (\textit{Learning With Errors}, o Aprendizaje con Errores), el juego del teléfono descompuesto controlado.
  \item \textbf{Capítulo 4:} Descubriremos por qué no podemos trabajar solo con números comunes y cómo los polinomios y los \textbf{módulos} nos permiten empaquetar terabytes de seguridad en el tamaño de una estampilla.
  \item \textbf{Capítulo 5:} Veremos el candado post-cuántico en acción: el estándar mundial \textbf{ML-KEM} (\textit{Kyber}), cómo Alicia y Roberto acuerdan una clave secreta a la vista de todos.
  \item \textbf{Capítulo 6:} Aprenderemos cómo funciona el sello post-cuántico: el estándar mundial \textbf{ML-DSA} (\textit{Dilithium}), el arte de firmar sin dejar pistas de tu clave secreta.
  \item \textbf{Capítulo 7:} Reflexionaremos sobre lo aprendido y te entregaremos un glosario amigable para que consultes cualquier término siempre que lo necesites.
\end{itemize}

\begin{intuicion}[Una invitación al asombro]
No leas este libro como quien memoriza fórmulas para un examen de medianoche. Míralo como una visita guiada por una galería de arte geométrico. Cada diagrama de TikZ y cada caja didáctica fueron dibujados para que tu intuición visual haga el trabajo pesado. ¡Que disfrutes el viaje!
\end{intuicion}

\chapter{La Pregunta Fundamental: ¿Cómo Guardar un Secreto con Matemáticas?}
\label{cap:el_arte_del_secreto}

\section{El dilema de la puerta trampa}

¿Alguna vez has intentado mezclar dos colores de pintura, por ejemplo amarillo y azul? 
En apenas dos segundos con un pincel obtienes un verde homogéneo. La operación directa es increíblemente sencilla. Pero ahora piensa en la operación inversa: te entregan un frasco lleno de pintura verde y te piden que separes cada molécula de amarillo y cada molécula de azul en dos vasos distintos. 

Es prácticamente imposible.

\begin{figure}[htbp]
\centering
\begin{tikzpicture}[scale=0.95]
  % Caja de entrada
  \node[draw=PQCMainBlue, fill=BoxBgBlue, rounded corners=3mm, line width=1pt, minimum width=2.8cm, minimum height=1.6cm, align=center] (in) at (0,0) {\textbf{Mensaje secreto}\\\small (ej. una clave o un voto)};
  
  % Caja de proceso
  \node[draw=PQCTeal, fill=BoxBgTeal, rounded corners=3mm, line width=1pt, minimum width=3.2cm, minimum height=1.6cm, align=center] (proc) at (5,0) {\textbf{Cálculo matemático}\\\small Fácil de calcular\\(Cerrar el candado)};
  
  % Caja de salida
  \node[draw=PQCPurple, fill=BoxBgPurple, rounded corners=3mm, line width=1pt, minimum width=2.8cm, minimum height=1.6cm, align=center] (out) at (10,0) {\textbf{Texto cifrado}\\\small Parece ruido al azar};
  
  % Flechas
  \draw[-{Stealth[length=3mm]}, line width=1.5pt, PQCMainBlue] (in) -- (proc) node[midway, above] {\small Entrada};
  \draw[-{Stealth[length=3mm]}, line width=1.5pt, PQCTeal] (proc) -- (out) node[midway, above] {\small Salida pública};
  
  % Flecha inversa imposible
  \draw[-{Stealth[length=3mm]}, line width=1.5pt, PQCCrimson, dashed] (out) to[bend left=35] node[midway, above] {\textbf{Invertir: Imposible sin la trampa}} (in);
  
  % Llave trampa
  \node[draw=PQCGreen, fill=BoxBgGreen, rounded corners=2mm, line width=1pt, below=0.8cm of proc, align=center] (key) {\textbf{Clave Privada (Trampa Secreta)}\\\small Permite regresar fácilmente};
  \draw[-{Stealth[length=2.5mm]}, line width=1.2pt, PQCGreen] (key) -- (proc);
\end{tikzpicture}
\caption{El principio de una función unidireccional con trampa (\textit{trapdoor function}).}
\label{fig:trampa}
\end{figure}

En matemáticas, a este tipo de procesos se les llama \textbf{funciones unidireccionales} (\textit{one-way functions}):
\begin{itemize}
  \item \textbf{Fáciles de calcular} en un sentido (calcular $f(x)$ dado $x$).
  \item \textbf{Prácticamente imposibles de invertir} para quien solo ve el resultado ($f(x)$), a menos que posea una información privilegiada llamada \textbf{puerta trampa} (\textit{trapdoor}).
\end{itemize}

\section{El problema de la escuela: La trampa de la linealidad}

En tus primeras semanas de álgebra lineal, aprendiste a resolver sistemas de ecuaciones como este:

\begin{align*}
2x + 3y &= 8 \\
5x + y &= 7
\end{align*}

¿Es difícil para una computadora resolver este sistema? En absoluto. El método de eliminación de Gauss o la regla de Cramer encuentran la solución exacta $(x=1, y=2)$ en una millonésima de segundo, incluso si el sistema tuviera $100$ ecuaciones con $100$ incógnitas.

\begin{ojo}[La linealidad es enemiga del secreto]
Si un candado criptográfico se basara únicamente en ecuaciones lineales ordinarias $A \cdot x = b$, cualquier atacante con conocimientos de álgebra escolar podría despejar el secreto $x$ en un abrir y cerrar de ojos. La precisión perfecta de los números permite al adversario cancelar variables con facilidad.
\end{ojo}

\section{La chispa genial: ¿Qué pasa si le sumamos un poco de ``ruido''?}

Imagina ahora que el maestro no te entrega las respuestas exactas, sino que a cada resultado le suma o le resta un pequeño número secreto, un ``error'' diminuto:

\begin{align*}
2x + 3y + e_1 &= 8.1 \\
5x + y + e_2 &= 6.9
\end{align*}
donde tú no conoces $e_1$ ni $e_2$, solo sabes que son números muy pequeños (quizás $-0.1, 0, 0.1$).

De pronto, todo el método de Gauss se derrumba:
\begin{itemize}
  \item Al multiplicar una ecuación por otra y restarlas, los pequeños errores ya no se cancelan: ¡se acumulan y explotan!
  \item Al cabo de unas cuantas operaciones, el resultado obtenido está completamente contaminado por el ruido.
\end{itemize}

\begin{metafora}[El juego del susurro en una fiesta ruidosa]
Imagina que un amigo te susurra una palabra en medio de un concierto de rock. Si no conoces de qué tema están hablando, escucharás un murmullo indescifrable. Pero si ambos comparten un contexto secreto previo, tu cerebro es capaz de filtrar el ruido ambiental y entender la palabra original a la perfección.
\end{metafora}

Esta idea tan simple y humana ---agregar un pequeño ruido deliberado que arruina la precisión del atacante pero que puede ser eliminado por quien tiene la clave correcta--- es el cimiento de toda la criptografía basada en retículos y de los estándares del siglo XXI.

\begin{resumencapitulo}
\begin{itemize}
  \item La criptografía moderna necesita operaciones fáciles de hacer pero casi imposibles de deshacer sin una clave secreta.
  \item Las ecuaciones lineales exactas son demasiado fáciles de resolver; no sirven por sí solas para ocultar secretos.
  \item Al agregar una pizca de ``ruido'' o error diminuto a las ecuaciones, los métodos clásicos de resolución (como Gauss) fallan por completo.
  \item Esta perturbación controlada nos llevará directamente al concepto de \textbf{Retículos} y al problema \textbf{LWE}.
\end{itemize}
\end{resumencapitulo}

\chapter{Una Cuadrícula Infinita en el Espacio: ¿Qué es un Retículo?}
\label{cap:reticulos_geometria}

\section{De la hoja cuadriculada al infinito}

Toma un cuaderno cuadriculado escolar. Cada intersección de dos líneas es un punto $(x, y)$ donde las coordenadas $x$ e $y$ son números enteros:
\[
(0,0),\; (1,0),\; (0,1),\; (2,3),\; (-1, 4),\dots
\]
Ese conjunto de puntos discretos y perfectamente espaciados que se extiende hacia el infinito se llama en matemáticas $\mathbb{Z}^2$.

Ahora imagina que deformamos esa cuadrícula: en lugar de caminar en pasos perpendiculares de longitud $1$ hacia la derecha y hacia arriba, elegimos dos flechas (vectores) cualesquiera en el plano, digamos:
\[
\mathbf{b}_1 = \begin{pmatrix} 2 \\ 1 \end{pmatrix}, \quad \mathbf{b}_2 = \begin{pmatrix} 1 \\ 3 \end{pmatrix}
\]

\begin{figure}[htbp]
\centering
\begin{tikzpicture}[scale=0.85]
  % Puntos del reticulo
  \foreach \x in {-2,...,3} {
    \foreach \y in {-2,...,2} {
      \coordinate (P) at ({\x*1.5 + \y*0.7}, {\x*0.4 + \y*1.8});
      \fill[PQCMainBlue!80!black] (P) circle (2.2pt);
    }
  }
  
  % Origen
  \coordinate (O) at (0,0);
  \node[below left, font=\small\bfseries] at (O) {$(0,0)$};
  
  % Vectores base b1 y b2
  \draw[-{Stealth[length=3mm]}, line width=1.8pt, PQCTeal] (O) -- (1.5, 0.4) node[midway, below right, font=\small\bfseries] {$\mathbf{b}_1$};
  \draw[-{Stealth[length=3mm]}, line width=1.8pt, PQCGreen] (O) -- (0.7, 1.8) node[midway, above left, font=\small\bfseries] {$\mathbf{b}_2$};
  
  % Combinacion de ejemplo
  \draw[-{Stealth[length=2.5mm]}, line width=1.2pt, PQCPurple, dashed] (O) -- ({2*1.5 + 1*0.7}, {2*0.4 + 1*1.8}) node[above right, font=\small\bfseries] {$2\mathbf{b}_1 + \mathbf{b}_2$};

\end{tikzpicture}
\caption{Un retículo en 2 dimensiones generado por los vectores $\mathbf{b}_1$ y $\mathbf{b}_2$. Cada punto es el resultado de sumar múltiplos enteros de las flechas base.}
\label{fig:reticulo_2d}
\end{figure}

Si sumamos o restamos estas dos flechas multiplicadas por números enteros (por ejemplo $3\mathbf{b}_1 - 2\mathbf{b}_2$), generamos una red infinita de puntos regulares en el plano.

\begin{intuicion}[Definición para primer semestre: ¿Qué es un Retículo?]
Un \textbf{retículo} (\textit{lattice} en inglés) es una cuadrícula infinita de puntos en el espacio donde:
\begin{enumerate}
  \item El origen $(0,0,\dots,0)$ siempre forma parte de la cuadrícula.
  \item Si tomas dos puntos cualquiera de la red y los sumas, caes exactamente en otro punto de la red.
  \item Si caminas de un punto al siguiente, siempre hay una distancia mínima positiva (los puntos están separados, no se amontonan de manera continua como el agua).
\end{enumerate}
\end{intuicion}

\section{Bases buenas y bases malas: El poder de la perspectiva}

Aquí viene el primer gran secreto de los retículos. Un mismo retículo (la misma constelación idéntica de puntos) puede ser generado por \textbf{diferentes pares de vectores base}.

\begin{itemize}
  \item \textbf{Una base ``buena'':} Está formada por vectores cortos y casi perpendiculares (a casi 90 grados). Es como tener una regla bien calibrada: es fácil saber en qué celda estás y a qué punto estás más cerca.
  \item \textbf{Una base ``mala'':} Está formada por vectores larguísimos, muy inclinados y casi paralelos entre sí. Aunque generan la misma red, visualmente parecen un laberinto enmarañado.
\end{itemize}

\begin{figure}[htbp]
\centering
\begin{tikzpicture}[scale=0.85]
  % Puntos
  \foreach \x in {-2,...,4} {
    \foreach \y in {-2,...,3} {
      \coordinate (P) at ({\x*1.3 + \y*0.3}, {\x*0.2 + \y*1.4});
      \fill[gray!60] (P) circle (1.8pt);
    }
  }
  
  \coordinate (O) at (0,0);
  
  % Base buena (verde)
  \draw[-{Stealth[length=2.5mm]}, line width=1.6pt, PQCGreen] (O) -- (1.3, 0.2) node[below, font=\footnotesize\bfseries] {$\mathbf{b}_1$ (Corta)};
  \draw[-{Stealth[length=2.5mm]}, line width=1.6pt, PQCGreen] (O) -- (0.3, 1.4) node[left, font=\footnotesize\bfseries] {$\mathbf{b}_2$ (Corta)};
  
  % Base mala (rojo)
  \draw[-{Stealth[length=2.5mm]}, line width=1.6pt, PQCCrimson] (O) -- ({3*1.3 + 2*0.3}, {3*0.2 + 2*1.4}) node[above right, font=\footnotesize\bfseries] {$\mathbf{v}_1$ (Larga y distorsionada)};
  \draw[-{Stealth[length=2.5mm]}, line width=1.6pt, PQCCrimson] (O) -- ({2*1.3 + 1*0.3}, {2*0.2 + 1*1.4}) node[midway, below right, font=\footnotesize\bfseries] {$\mathbf{v}_2$};

\end{tikzpicture}
\caption{La misma cuadrícula vista a través de dos bases distintas: la base buena (verde) es clara y ordenada; la base mala (roja) es confusa y distorsionada.}
\label{fig:base_buena_mala}
\end{figure}

\begin{metafora}[La cuadrícula de una ciudad]
Imagina una ciudad como Manhattan, con avenidas y calles perpendiculares (base buena). Si te dicen ``avanza 3 cuadras al norte y 2 al este'', sabes exactamente dónde estás y cuál es la casa más cercana. 

Ahora imagina que te dan dos direcciones raras: una avenida diagonal que avanza 100 kilómetros hacia el noreste y otra que avanza 99 kilómetros casi en la misma dirección (base mala). Para llegar a la tienda de la esquina tendrías que sumar y restar números gigantescos con precisión astronómica. Cualquier pequeño error te dejaría en otra provincia.
\end{metafora}

\section{El problema geométrico central: CVP (El vector más cercano)}

Ahora presentemos el problema que le quita el sueño a las supercomputadoras: el \textbf{Problema del Vector Más Cercano} (\textit{Closest Vector Problem} o CVP).

\begin{intuicion}[El juego del CVP]
\textbf{Planteamiento:} Imagina que te entregan un punto en el espacio $\mathbf{t}$ que no pertenece al retículo (cayó en el ``espacio vacío'' entre los puntos de la cuadrícula, por culpa de un pequeño error). 

\textbf{El desafío:} ¿Cuál es el punto del retículo más cercano a $\mathbf{t}$?
\end{intuicion}

\begin{figure}[htbp]
\centering
\begin{tikzpicture}[scale=0.95]
  % Puntos
  \foreach \x in {-1,...,3} {
    \foreach \y in {-1,...,2} {
      \coordinate (P) at ({\x*1.8 + \y*0.4}, {\x*0.3 + \y*1.6});
      \fill[PQCMainBlue] (P) circle (2.2pt);
    }
  }
  
  % Punto objetivo t (fuera de la malla)
  \coordinate (T) at ({1*1.8 + 1*0.4 + 0.35}, {1*0.3 + 1*1.6 + 0.45});
  \fill[PQCCrimson] (T) circle (3pt) node[above right, font=\small\bfseries] {Punto objetivo $\mathbf{t}$ (con ruido)};
  
  % Punto mas cercano
  \coordinate (Closest) at ({1*1.8 + 1*0.4}, {1*0.3 + 1*1.6});
  \draw[PQCGreen, line width=1.5pt] (Closest) circle (6pt);
  \draw[-{Stealth[length=2.5mm]}, line width=1.2pt, PQCAmber] (Closest) -- (T) node[midway, below right, font=\footnotesize\bfseries] {ruido $\mathbf{e}$};
  \node[below left, font=\footnotesize\bfseries, PQCGreen] at (Closest) {¡Punto más cercano!};
\end{tikzpicture}
\caption{El problema CVP: Dado un punto $\mathbf{t}$ que fue desplazado por un ruido $\mathbf{e}$, el objetivo es encontrar cuál punto azul era el original.}
\label{fig:cvp}
\end{figure}

\begin{itemize}
  \item \textbf{Si tienes la base buena:} Miras tus coordenadas, redondeas al entero más cercano, y en un instante descubres el punto exacto.
  \item \textbf{Si solo tienes la base mala:} No sabes hacia dónde redondear. En dimensiones altas (por ejemplo $n = 500$ o $n = 768$ dimensiones), hay tantas posibilidades de búsqueda que ni siquiera todas las computadoras del planeta trabajando juntas durante siglos podrían encontrar el punto correcto.
\end{itemize}

¡Eureka! Ya tenemos nuestra primera trampa criptográfica:
\begin{itemize}
  \item \textbf{Tu Clave Privada} es la \textbf{base buena} (corta y ordenada).
  \item \textbf{Tu Clave Pública} es la \textbf{base mala} (larga y desordenada).
  \item Cualquiera puede usar tu clave pública para sumar vectores y agregarle un poco de ruido (cifrar).
  \item Solo tú, con tu base buena, puedes quitar el ruido y encontrar el mensaje original (descifrar).
\end{itemize}

\begin{resumencapitulo}
\begin{itemize}
  \item Un retículo es una cuadrícula infinita y periódica de puntos generada por combinaciones lineales enteras de vectores.
  \item Un mismo retículo tiene infinitas bases: bases ``buenas'' (vectores cortos, fáciles de manejar) y bases ``malas'' (vectores largos y casi paralelos).
  \item El Problema del Vector Más Cercano (CVP) consiste en hallar el punto del retículo más próximo a un punto ruidoso.
  \item En cientos de dimensiones, CVP es extremadamente fácil si conoces la base buena, pero computacionalmente imposible si solo tienes la base mala.
\end{itemize}
\end{resumencapitulo}

\chapter{El Secreto con Ruido: Aprendiendo con Errores (LWE)}
\label{cap:lwe_ruido}

\section{La aritmética del reloj: Números que dan vueltas}

Antes de comprender cómo funciona el problema más famoso de la criptografía moderna, necesitamos recordar una herramienta matemática que utilizas todos los días sin darte cuenta: \textbf{el reloj}.

Si tu reloj marca las $10$ de la mañana y tu clase dura $4$ horas, no dices que saldrás a las ``$14$ de la mañana'', sino a las $2$. ¿Qué hiciste mentalmente? Sumaste $10 + 4 = 14$, le restaste el ciclo completo de $12$, y te quedaste con el residuo:
\[
10 + 4 \equiv 2 \pmod{12}
\]

\begin{figure}[htbp]
\centering
\begin{tikzpicture}[scale=0.95]
  % Circulo
  \draw[line width=1.5pt, PQCMainBlue] (0,0) circle (2.2cm);
  
  % Horas 0 a 10 (Z_11)
  \foreach \i in {0,...,10} {
    \coordinate (H\i) at ({90 - \i*360/11}:2.2cm);
    \fill[PQCMainBlue] (H\i) circle (3pt);
    \node[font=\small\sffamily\bfseries] at ({90 - \i*360/11}:2.65cm) {\i};
  }
  
  % Flecha que ilustra dar la vuelta
  \draw[-{Stealth[length=3mm]}, line width=1.5pt, PQCTeal] (90:1.6cm) arc (90:-40:1.6cm);
  \node[font=\footnotesize\sffamily, align=center, PQCTeal] at (0.4, 0.5) {$8 + 5 = 13 \equiv 2$\\(en ciclo $q=11$)};
\end{tikzpicture}
\caption{La aritmética modular $\mathbb{Z}_{11}$: una rueda cerrada donde los números nunca escapan al infinito.}
\label{fig:reloj_modular}
\end{figure}

En criptografía, en lugar de trabajar con la recta real infinita $\mathbb{R}$ (donde los números con decimales crecen y consumen memoria infinita), fijamos un número entero positivo $q$ llamado \textbf{módulo}. 

El universo donde vivimos se llama $\mathbb{Z}_q = \{0, 1, 2, \dots, q-1\}$. Cada vez que una suma o multiplicación sobrepasa $q$, simplemente nos quedamos con el residuo de la división entre $q$.

\section{El problema LWE: El invento de Oded Regev}

En el año 2005, el matemático y científico de la computación Oded Regev formuló un problema computacional brillante llamado \textbf{LWE} (\textit{Learning With Errors}, o Aprendizaje con Errores). 

Su idea es tan elegante que cabe en una conversación de café:

\begin{intuicion}[¿Cómo se juega a LWE?]
\begin{enumerate}
  \item \textbf{El Secreto:} Imagina que tienes un vector secreto $\mathbf{s}$ compuesto por números enteros pequeños (por ejemplo $\mathbf{s} = (s_1, s_2)$).
  \item \textbf{Las Preguntas Públicas:} Alguien elige al azar una fila de números $\mathbf{a} = (a_1, a_2)$ a la vista de todo el mundo.
  \item \textbf{El Cálculo Limpio:} Haces el producto punto habitual:
  \[
  \mathbf{a} \cdot \mathbf{s} = a_1 s_1 + a_2 s_2 \pmod q
  \]
  \item \textbf{El Ruido Mágico:} En lugar de publicar ese resultado limpio, le sumas un error diminuto secreto $e$ (elegido al azar, por ejemplo de $\{-1, 0, 1\}$):
  \[
  b = \mathbf{a} \cdot \mathbf{s} + e \pmod q
  \]
\end{enumerate}
\end{intuicion}

Tú publicas los pares $(\mathbf{a}, b)$. Todo el mundo conoce las preguntas $\mathbf{a}$ y los resultados ruidosos $b$. 

\begin{ojo}[El desafío imposible para el atacante]
Dada una lista pública gigante de parejas $(\mathbf{a}_1, b_1), (\mathbf{a}_2, b_2), \dots, (\mathbf{a}_m, b_m)$, ¿puede un atacante descubrir el vector secreto $\mathbf{s}$?

\textbf{Respuesta:} ¡No! Si no hubiera error ($e=0$), el sistema se resolvería con el método de Gauss en milisegundos. Pero al existir el pequeño error desconocido $e$, el problema de hallar $\mathbf{s}$ es computacionalmente equivalente a resolver el problema del vector más cercano (CVP) en un retículo de alta dimensión.
\end{ojo}

\section{Un ejemplo numérico paso a paso}

Hagamos juntos un ejercicio real con números pequeños para ver la magia con tus propios ojos.

\begin{ejemplonumerico}[LWE con módulo $q=13$]
Elijamos un módulo pequeño $q = 13$.
Supongamos que nuestro secreto es un solo número: $s = 4$.

Generemos tres muestras públicas:
\begin{enumerate}
  \item \textbf{Muestra 1:} Elegimos al azar $a_1 = 3$. 
  Calculamos $3 \times 4 = 12$. Agregamos un pequeño ruido $e_1 = +1$:
  \[
  b_1 = 12 + 1 = 13 \equiv 0 \pmod{13}
  \]
  Publicamos el par $(3, 0)$.
  
  \item \textbf{Muestra 2:} Elegimos al azar $a_2 = 5$. 
  Calculamos $5 \times 4 = 20 \equiv 7 \pmod{13}$. Agregamos un ruido $e_2 = -1$:
  \[
  b_2 = 7 - 1 = 6 \pmod{13}
  \]
  Publicamos el par $(5, 6)$.
  
  \item \textbf{Muestra 3:} Elegimos al azar $a_3 = 2$. 
  Calculamos $2 \times 4 = 8$. Ruido $e_3 = 0$:
  \[
  b_3 = 8 + 0 = 8 \pmod{13}
  \]
  Publicamos el par $(2, 8)$.
\end{enumerate}

El público solo ve:
\[
(a_1=3, b_1=0), \quad (a_2=5, b_2=6), \quad (a_3=2, b_3=8)
\]
Para un atacante, los números $b$ parecen absolutamente aleatorios. No puede simplemente ``despejar'' $s$ porque no sabe qué porción de cada $b$ es información pura y qué porción es ruido.
\end{ejemplonumerico}

\section{¿Por qué una computadora cuántica no puede romper LWE?}

Las computadoras cuánticas son máquinas fascinantes, pero no son varitas mágicas capaces de probar todas las combinaciones al instante. 

El motivo por el cual una computadora cuántica puede resolver ciertos problemas tradicionales es que esos problemas antiguos poseen una \textbf{estructura periódica regular} (como una onda de radio limpia con una frecuencia exacta). El algoritmo cuántico de Shor es una especie de ``sintonizador de radio cuántico'' ultrasensible que detecta periodos regulares a la perfección.

\begin{intuicion}[La belleza del caos en LWE]
En el problema LWE, el ruido rompe cualquier patrón periódico. No hay una onda limpia que sintonizar: cada ecuación fue perturbada en una dirección diferente e impredecible. Por eso, ni las supercomputadoras clásicas más potentes ni las futuras computadoras cuánticas universales conocen un algoritmo eficiente para quebrar LWE en dimensiones altas.
\end{intuicion}

\begin{resumencapitulo}
\begin{itemize}
  \item La aritmética modular $\mathbb{Z}_q$ confina todos los cálculos a un anillo finito de números como las horas de un reloj.
  \item LWE formula ecuaciones lineales simples perturbadas por un pequeño error secreto: $\mathbf{b} = \mathbf{A}\mathbf{s} + \mathbf{e} \pmod q$.
  \item El ruido destruye la posibilidad de aplicar eliminación gaussiana, convirtiendo el problema en una búsqueda geométrica en un retículo.
  \item La ausencia de periodicidad regular protege a LWE tanto de atacantes convencionales como de computadoras cuánticas.
\end{itemize}
\end{resumencapitulo}

\chapter{De Números Sueltos a Polinomios: El Salto a Module-LWE}
\label{cap:polinomios_modulos}

\section{El problema de la maleta pesada}

Hasta ahora, la idea de LWE parece perfecta: multiplicamos matrices por vectores, sumamos un pequeño error y tenemos un candado impenetrable. 

Sin embargo, cuando los ingenieros intentaron llevar LWE puro a la práctica en teléfonos celulares y tarjetas inteligentes, se toparon con una pared de concreto: \textbf{el tamaño}.

Para que LWE clásico sea seguro contra computadoras de alta potencia, la matriz pública $\mathbf{A}$ debe tener dimensiones gigantescas (por ejemplo, $1000 \times 1000$ números enteros). Almacenar y transmitir esa matriz por internet requeriría varios megabytes en cada clic. Imagina enviar un mensaje por WhatsApp o emitir un voto en una elección digital y tener que esperar segundos solo para transferir las claves.

\begin{metafora}[El diccionario frente a la biblioteca]
LWE estándar es como cargar una enciclopedia de $20$ tomos en la mochila cada vez que vas a la universidad. Los matemáticos necesitaban una forma de empaquetar toda esa información en una tarjeta de bolsillo elegante y liviana.
\end{metafora}

¿Cuál fue la solución mágica? Un concepto que conoces desde el colegio: \textbf{los polinomios}.

\section{El carrusel de los polinomios: \texorpdfstring{$x^n + 1 = 0$}{x\textasciicircum n + 1 = 0}}

En la secundaria aprendiste a sumar y multiplicar polinomios sencillos como $(2x + 1)(x + 3)$. 

En la criptografía post-cuántica moderna trabajamos en un anillo especial llamado:
\[
R_q = \mathbb{Z}_q[x]/(x^n + 1)
\]
El nombre puede sonar intimidante para un estudiante de primer semestre, pero la idea física es tan sencilla como un carrusel de feria:

\begin{intuicion}[La regla del carrusel]
Un polinomio en $R_q$ solo tiene términos de grado menor a $n$:
\[
P(x) = c_0 + c_1 x + c_2 x^2 + \dots + c_{n-1} x^{n-1}
\]
donde cada coeficiente $c_i$ vive en el reloj $\mathbb{Z}_q$.

¿Qué pasa si al multiplicar dos polinomios aparece una potencia alta como $x^n$? Aplicamos la regla mágica de reducción:
\[
x^n \equiv -1
\]
Cada vez que una potencia alcanza $x^n$, ``da la vuelta'' al carrusel, reaparece como una potencia baja y cambia de signo:
\[
x^n = -1, \quad x^{n+1} = -x, \quad x^{n+2} = -x^2, \quad \dots
\]
\end{intuicion}

\begin{figure}[htbp]
\centering
\begin{tikzpicture}[scale=1.0]
  % Dibujo de carrusel con 4 potencias (n=4)
  \coordinate (O) at (0,0);
  \draw[line width=1.5pt, PQCTeal, dashed] (O) circle (2.0cm);
  
  \node[draw=PQCMainBlue, fill=BoxBgBlue, circle, font=\small\bfseries] (x0) at (0,2.0) {$x^0=1$};
  \node[draw=PQCMainBlue, fill=BoxBgBlue, circle, font=\small\bfseries] (x1) at (2.0,0) {$x^1$};
  \node[draw=PQCMainBlue, fill=BoxBgBlue, circle, font=\small\bfseries] (x2) at (0,-2.0) {$x^2$};
  \node[draw=PQCMainBlue, fill=BoxBgBlue, circle, font=\small\bfseries] (x3) at (-2.0,0) {$x^3$};
  
  % Flechas
  \draw[-{Stealth[length=2.5mm]}, line width=1.4pt, PQCMainBlue] (x0) to[bend left=20] (x1);
  \draw[-{Stealth[length=2.5mm]}, line width=1.4pt, PQCMainBlue] (x1) to[bend left=20] (x2);
  \draw[-{Stealth[length=2.5mm]}, line width=1.4pt, PQCMainBlue] (x2) to[bend left=20] (x3);
  
  % El salto de cambio de signo
  \draw[-{Stealth[length=3mm]}, line width=1.8pt, PQCCrimson] (x3) to[bend left=30] node[midway, above left, font=\footnotesize\bfseries, align=center] {¡Multiplicar por $x$\\\textbf{da $-1$}!} (x0);
\end{tikzpicture}
\caption{El carrusel de potencias en el anillo con $n=4$: cuando llegamos a $x^4$, rebotamos al inicio con signo cambiado ($-1$). Ningún polinomio se desborda en grado.}
\label{fig:carrusel_polinomios}
\end{figure}

\section{¿Qué es un ``Módulo'' en palabras sencillas?}

En álgebra lineal de primer semestre, un vector es una columna de números comunes y corrientes:
\[
\mathbf{v} = \begin{pmatrix} 3 \\ -5 \\ 8 \end{pmatrix} \in \mathbb{R}^3
\]
Un \textbf{módulo} (en el sentido de retículos modulares) es exactamente la misma estructura, pero con un detalle fabuloso: \textbf{en lugar de poner números en cada coordenada, ponemos polinomios completos}.

\[
\mathbf{v} = \begin{pmatrix} P_1(x) \\ P_2(x) \\ P_3(x) \end{pmatrix} \in R_q^3
\]

Por ejemplo, en el estándar de seguridad \textbf{ML-KEM-768}, tenemos:
\begin{itemize}
  \item Cada polinomio tiene grado $n = 256$ (es decir, almacena $256$ coeficientes).
  \item El vector tiene dimensión $k = 3$ (contiene $3$ polinomios).
  \item ¡En total estamos manejando $3 \times 256 = 768$ coeficientes matemáticos a la vez!
\end{itemize}

\section{Module-LWE (M-LWE): La joya de la corona}

Ahora combinamos todo lo aprendido:
\begin{equation}
\mathbf{b} = \mathbf{A} \cdot \mathbf{s} + \mathbf{e} \pmod q
\end{equation}
¡Pero ahora $\mathbf{A}$ no es una matriz colosal de millones de números! Es una matriz minúscula de $3 \times 3$ polinomios.

\begin{figure}[htbp]
\centering
\begin{tikzpicture}[scale=0.95]
  % Matriz A
  \node[draw=PQCMainBlue, fill=BoxBgBlue, line width=1pt, rounded corners=2mm, minimum width=2.4cm, minimum height=1.8cm, align=center] (A) at (0,0) {$\mathbf{A}$\\\footnotesize Matriz $3 \times 3$\\\scriptsize de polinomios};
  
  \node[font=\large\bfseries] at (1.6,0) {$\cdot$};
  
  % Vector s
  \node[draw=PQCPurple, fill=BoxBgPurple, line width=1pt, rounded corners=2mm, minimum width=1.5cm, minimum height=1.8cm, align=center] (s) at (2.7,0) {$\mathbf{s}$\\\footnotesize Secreto\\\scriptsize (3 polinomios)};
  
  \node[font=\large\bfseries] at (3.8,0) {$+$};
  
  % Vector e
  \node[draw=PQCAmber, fill=BoxBgAmber, line width=1pt, rounded corners=2mm, minimum width=1.5cm, minimum height=1.8cm, align=center] (e) at (4.9,0) {$\mathbf{e}$\\\footnotesize Ruido\\\scriptsize (3 polinomios)};
  
  \node[font=\large\bfseries] at (6.0,0) {$=$};
  
  % Vector b
  \node[draw=PQCTeal, fill=BoxBgTeal, line width=1pt, rounded corners=2mm, minimum width=1.5cm, minimum height=1.8cm, align=center] (b) at (7.1,0) {$\mathbf{b}$\\\footnotesize Clave pública\\\scriptsize (3 polinomios)};
\end{tikzpicture}
\caption{La ecuación fundamental de Module-LWE: dimensiones pequeñas en álgebra matricial, pero enorme riqueza geométrica interna gracias a los polinomios.}
\label{fig:mlwe_arquitectura}
\end{figure}

\begin{intuicion}[¿Por qué Module-LWE es la mejor opción del mundo?]
\begin{enumerate}
  \item \textbf{Velocidad ultrasónica:} Los polinomios se pueden multiplicar casi instantáneamente usando el equivalente algebraico de la Transformada Rápida de Fourier (llamada NTT, \textit{Number Theoretic Transform}).
  \item \textbf{Tamaño de bolsillo:} Las claves públicas miden apenas unos mil bytes (caben cómodamente en un solo paquete de red de internet).
  \item \textbf{Seguridad garantizada:} Romper este esquema obligaría al adversario a resolver problemas geométricos insolubles en retículos modulares de más de $700$ dimensiones.
\end{enumerate}
\end{intuicion}

\begin{resumencapitulo}
\begin{itemize}
  \item LWE estándar requiere matrices demasiado pesadas para aplicaciones cotidianas.
  \item El anillo de polinomios $R_q = \mathbb{Z}_q[x]/(x^n+1)$ permite comprimir la información mediante un carrusel cíclico de potencias ($x^n = -1$).
  \item Un módulo es un espacio vectorial donde cada coordenada es un polinomio en lugar de un escalar.
  \item \textbf{Module-LWE (M-LWE)} combina matrices pequeñas ($3 \times 3$) con polinomios ($n=256$), logrando claves ultracompactas, cómputo ultrarrápido y máxima resistencia cuántica.
\end{itemize}
\end{resumencapitulo}

\chapter{El Candado Post-Cuántico: ¿Cómo Funciona ML-KEM?}
\label{cap:mlkem_candado}

\section{¿Qué es un KEM y por qué lo usamos?}

En las películas de espías, cuando alguien quiere enviar un mensaje secreto suele recurrir a una metáfora simple: cerrar una caja fuerte con un candado cuya llave solo tiene el destinatario.

En la informática moderna, los algoritmos basados en retículos son increíblemente potentes, pero no se usan directamente para cifrar archivos gigantescos de gigabytes (como videos o bases de datos completas). En su lugar, se utilizan para lo que los matemáticos llaman un \textbf{KEM} (\textit{Key Encapsulation Mechanism}, o Mecanismo de Encapsulamiento de Claves).

\begin{metafora}[La caja fuerte dentro del camión blindado]
Imagina que quieres enviar un camión lleno de documentos confidenciales (por ejemplo, millones de votos ciudadanos). No necesitas construir un camión entero de retículos post-cuánticos.
\begin{enumerate}
  \item Creas una pequeña llave simétrica ultrarrápida (como una llave AES-256).
  \item Guardas esa diminuta llave en una cajita fuerte blindada con retículos post-cuánticos (\textbf{ML-KEM}).
  \item El destinatario abre la cajita con su clave privada, extrae la llave y descifra el camión en un pestañeo.
\end{enumerate}
\end{metafora}

El estándar mundial definitivo, oficializado por el NIST en 2024 bajo el nombre \textbf{FIPS 203}, es \textbf{ML-KEM} (originalmente conocido como \textit{CRYSTALS-Kyber}). Sus parámetros de referencia son $n=256$, $q=3329$ y $k=3$, lo que da un nivel de seguridad equivalente a AES-192.

\section{La danza de tres pasos: Generar, Encapsular y Desencapsular}

El funcionamiento de ML-KEM consta de tres actos matemáticos ejecutados entre dos personas: \textbf{Roberto} (el receptor, quien crea la cerradura) y \textbf{Alicia} (la emisora, quien quiere mandar el secreto).

\begin{figure}[htbp]
\centering
\begin{tikzpicture}[scale=0.92]
  % Roberto y Alicia
  \node[draw=PQCMainBlue, fill=BoxBgBlue, rounded corners=2mm, line width=1pt, minimum width=3.2cm, minimum height=1.2cm, align=center] (R) at (0,4) {\textbf{Roberto}\\\small (Receptor / Autoridad)};
  \node[draw=PQCPurple, fill=BoxBgPurple, rounded corners=2mm, line width=1pt, minimum width=3.2cm, minimum height=1.2cm, align=center] (A) at (8,4) {\textbf{Alicia}\\\small (Emisora / Votante)};
  
  % Paso 1: Clave publica
  \draw[-{Stealth[length=3mm]}, line width=1.5pt, PQCMainBlue] (1.8,3.2) -- (6.2,3.2) node[midway, above, font=\small\sffamily] {1. Clave pública $(\mathbf{A}, \mathbf{b})$};
  
  % Paso 2: Texto cifrado
  \draw[-{Stealth[length=3mm]}, line width=1.5pt, PQCPurple] (6.2,1.8) -- (1.8,1.8) node[midway, above, font=\small\sffamily] {2. Boleta cifrada $c = (\mathbf{u}, v)$};
  
  % Paso 3: Secreto compartido
  \node[draw=PQCGreen, fill=BoxBgGreen, rounded corners=2mm, line width=1pt, align=center, font=\small\sffamily] at (0,0.4) {3. Desencapsula\\y recupera $K$};
  \node[draw=PQCGreen, fill=BoxBgGreen, rounded corners=2mm, line width=1pt, align=center, font=\small\sffamily] at (8,0.4) {Usa la misma\\clave simétrica $K$};
  
\end{tikzpicture}
\caption{Protocolo general de encapsulamiento con ML-KEM.}
\label{fig:danza_kem}
\end{figure}

\subsection*{Paso 1: Roberto crea su cerradura (Generación de claves)}
Roberto utiliza Module-LWE:
\begin{enumerate}
  \item Elige una matriz pública $\mathbf{A}$ de polinomios en $R_q^{k \times k}$.
  \item Escoge en secreto dos vectores de polinomios con coeficientes diminutos: el secreto $\mathbf{s}$ y el error $\mathbf{e}$.
  \item Calcula:
  \[
  \mathbf{b} = \mathbf{A} \cdot \mathbf{s} + \mathbf{e} \pmod q
  \]
  \item \textbf{Clave Pública (para todo el mundo):} $(\mathbf{A}, \mathbf{b})$.
  \item \textbf{Clave Privada (bajo siete llaves):} $\mathbf{s}$.
\end{enumerate}

\begin{ojo}[¿Cuánto mide la clave pública real?]
En ML-KEM-768, la clave pública completa ocupa apenas \textbf{1184 bytes} (aproximadamente 1 kilobyte). Para comparar, un archivo de música en baja calidad ocupa casi 1000 veces más. Esto la hace perfectamente transmisible por cualquier red de internet moderna, desde una conexión de banda ancha hasta un enlace 4G.
\end{ojo}

\subsection*{Paso 2: Alicia encapsula el secreto}
Alicia quiere enviarle a Roberto una clave secreta representada por un mensaje de bits $m \in \{0, 1\}^{256}$.
\begin{enumerate}
  \item Alicia elige su propio vector secreto aleatorio $\mathbf{r}$ y dos pequeños ruidos $\mathbf{e}_1, e_2$.
  \item Calcula dos paquetes matemáticos:
  \[
  \mathbf{u} = \mathbf{A}^T \mathbf{r} + \mathbf{e}_1 \pmod q
  \]
  \[
  v = \mathbf{b}^T \mathbf{r} + e_2 + \text{Codificar}(m) \pmod q
  \]
  \item Envía el par $c = (\mathbf{u}, v)$ a través de internet.
\end{enumerate}

La \textbf{intuición geométrica} del encapsulado es fascinante: Alicia está construyendo su propio ``par de claves efímero'' usando la clave pública de Roberto como punto de referencia, y escondiendo el secreto $m$ entre el ruido generado por ese par temporal. Sin la clave privada $\mathbf{s}$ de Roberto, el texto cifrado $(\mathbf{u}, v)$ es indistinguible de números completamente aleatorios.

\subsection*{Paso 3: Roberto abre el candado (La magia de la cancelación)}
Cualquier espía que intercepte $(\mathbf{u}, v)$ ve un laberinto de números incomprensibles. Pero Roberto tiene su llave mágica $\mathbf{s}$. 

Observa con atención el álgebra que aprende cualquier estudiante de primer semestre: Roberto multiplica $\mathbf{u}$ por su clave secreta $\mathbf{s}$:
\[
\mathbf{s}^T \mathbf{u} = \mathbf{s}^T (\mathbf{A}^T \mathbf{r} + \mathbf{e}_1) = \mathbf{s}^T \mathbf{A}^T \mathbf{r} + \mathbf{s}^T \mathbf{e}_1
\]

Ahora Roberto resta este resultado del valor $v$:
\[
v - \mathbf{s}^T \mathbf{u} = \Big(\mathbf{b}^T \mathbf{r} + e_2 + \text{Codificar}(m)\Big) - \Big(\mathbf{s}^T \mathbf{A}^T \mathbf{r} + \mathbf{s}^T \mathbf{e}_1\Big)
\]

Recordemos que la clave pública de Roberto era $\mathbf{b} = \mathbf{As} + \mathbf{e}$, por lo que $\mathbf{b}^T \mathbf{r} = \mathbf{s}^T \mathbf{A}^T \mathbf{r} + \mathbf{e}^T \mathbf{r}$. Al sustituir:
\[
v - \mathbf{s}^T \mathbf{u} = \text{Codificar}(m) + \underbrace{\big(\mathbf{e}^T \mathbf{r} + e_2 - \mathbf{s}^T \mathbf{e}_1\big)}_{\text{¡Un único residuo de ruido diminuto!}}
\]

\begin{intuicion}[¡Las montañas se cancelan, solo queda un susurro!]
Los términos gigantescos $(\mathbf{s}^T \mathbf{A}^T \mathbf{r})$ que ocultaban el secreto se restan a sí mismos y desaparecen por completo. Lo único que queda en manos de Roberto es el mensaje codificado más un pequeño puñado de ruido residual.
\end{intuicion}

\section{¿Cómo se elimina el ruido final? El truco del redondeo}

¿Cómo sabe Roberto si un bit original era un $0$ o un $1$ si todavía le queda ese pequeño residuo de ruido?

La solución es hermosamente geométrica:
\begin{itemize}
  \item Si Alicia quiere mandar un $0$, lo codifica en el punto \textbf{$0$} del reloj.
  \item Si Alicia quiere mandar un $1$, lo codifica exactamente en el punto opuesto del reloj, a mitad de camino: \textbf{$\lfloor q/2 \rceil$}.
\end{itemize}

Concretamente, la función de codificación trabaja así:
\[
\text{Codificar}(0) = 0 \qquad \text{Codificar}(1) = \left\lfloor \frac{q}{2} \right\rceil = 1665 \text{ (para } q = 3329\text{)}
\]

\begin{figure}[htbp]
\centering
\begin{tikzpicture}[scale=0.95]
  % Eje de 0 a q
  \draw[line width=1.5pt, gray!70] (0,0) -- (9,0);
  
  % Puntos clave
  \fill[PQCMainBlue] (0,0) circle (3pt) node[below=2mm, font=\small\bfseries] {$0$ (Bit 0)};
  \fill[PQCPurple] (4.5,0) circle (3pt) node[below=2mm, font=\small\bfseries] {$\lfloor q/2 \rceil = 1665$ (Bit 1)};
  \fill[PQCMainBlue] (9,0) circle (3pt) node[below=2mm, font=\small\bfseries] {$q=3329 \equiv 0$};
  
  % Zona de ruido alrededor de 0
  \fill[PQCMainBlue!20, rounded corners=2mm] (-0.8, 0.4) rectangle (1.2, 0.9);
  \node[font=\footnotesize\sffamily, PQCMainBlue] at (0.2, 0.65) {Zona bit 0};
  
  % Zona de ruido alrededor de q/2
  \fill[PQCPurple!20, rounded corners=2mm] (3.5, 0.4) rectangle (5.5, 0.9);
  \node[font=\footnotesize\sffamily, PQCPurple] at (4.5, 0.65) {Zona bit 1};
  
  % Umbrales de decision
  \draw[dashed, line width=1.2pt, PQCAmber] (2.25, -0.6) -- (2.25, 1.2) node[above, font=\footnotesize\sffamily] {Frontera $q/4$};
  \draw[dashed, line width=1.2pt, PQCAmber] (6.75, -0.6) -- (6.75, 1.2) node[above, font=\footnotesize\sffamily] {Frontera $3q/4$};
  
  % Ruido simulado
  \draw[-{Stealth[length=2mm]}, PQCCrimson, line width=1pt] (4.5, -0.3) -- (5.0, -0.3) node[right, font=\scriptsize, PQCCrimson] {Ruido};
  
\end{tikzpicture}
\caption{Decodificación por redondeo: el ruido desplaza ligeramente el punto, pero jamás cruza las fronteras de decisión en $q/4$ y $3q/4$.}
\label{fig:redondeo_kem}
\end{figure}

Como los ruidos elegidos son matemáticamente muy pequeños comparados con la distancia $q/4$, cuando Roberto calcula $v - \mathbf{s}^T \mathbf{u}$:
\begin{itemize}
  \item Si el valor cae en la zona $[0, q/4)$ o $(3q/4, q)$, Roberto concluye con certeza matemática: \textbf{¡El bit era un $0$!}
  \item Si el valor cae en la zona $(q/4, 3q/4)$, Roberto concluye: \textbf{¡El bit era un $1$!}
\end{itemize}

El ruido se esfuma, y el secreto queda recuperado a la perfección.

\section{¿Qué ocurre si un espía intercepta el mensaje?}

Una pregunta natural es: ¿qué ve exactamente un espía que captura el paquete $c = (\mathbf{u}, v)$ transmitido por internet?

\begin{figure}[htbp]
\centering
\begin{tikzpicture}[scale=0.9]
  % Alicia (izquierda)
  \node[draw=PQCPurple, fill=BoxBgPurple, rounded corners=2mm, line width=1pt, align=center, minimum width=2.4cm, minimum height=1.0cm] (A) at (0,0) {\textbf{Alicia}};
  
  % Espia (centro arriba)
  \node[draw=PQCCrimson, fill=red!8, rounded corners=2mm, line width=1pt, align=center, minimum width=2.8cm, minimum height=1.0cm] (E) at (5,1.5) {\textbf{Espía}\\{\scriptsize ve: números aleatorios}};
  
  % Roberto (derecha)
  \node[draw=PQCMainBlue, fill=BoxBgBlue, rounded corners=2mm, line width=1pt, align=center, minimum width=2.4cm, minimum height=1.0cm] (R) at (10,0) {\textbf{Roberto}};
  
  % Canal publico
  \draw[line width=2pt, PQCTeal] (A) -- node[below, font=\small\sffamily, align=center] {Canal público\\$(\mathbf{A}, c=(\mathbf{u},v))$} (R);
  
  % Espia captura
  \draw[dashed, line width=1.2pt, PQCCrimson, -{Stealth[length=2.5mm]}] (5, 0.35) -- (E);
  
  % Resultado del espia
  \node[font=\footnotesize\itshape, PQCCrimson, align=center] at (5, 3.0) {``Sé que $\mathbf{u}$ y $v$ existen,\\pero no puedo encontrar $m$ sin $\mathbf{s}$''};
  
\end{tikzpicture}
\caption{El espía intercepta $c = (\mathbf{u}, v)$, pero sin la clave privada $\mathbf{s}$ el cifrado es computacionalmente indistinguible de ruido.}
\label{fig:espia_kem}
\end{figure}

El espía solo ve:
\begin{enumerate}
  \item La matriz $\mathbf{A}$ (pública desde el inicio, no sirve sola).
  \item El vector $\mathbf{u}$ y el escalar $v$: parecen distribuidos uniformemente en $\mathbb{Z}_q$, es decir, \textbf{idénticos a ruido aleatorio}.
\end{enumerate}

Matemáticamente se demuestra que distinguir un texto cifrado ML-KEM de una cadena de bits puramente aleatoria es tan difícil como resolver el problema Module-LWE en dimensión 768, lo cual se estima computacionalmente imposible para los próximos siglos, incluso usando computadoras cuánticas.

\section{Rendimiento real: ¿Cuánto tarda ML-KEM?}

Una preocupación natural cuando se menciona ``matemática post-cuántica'' es que todo sea enormemente lento. Los números reales desmitifican esta idea de manera rotunda:

\begin{figure}[htbp]
\centering
\begin{tikzpicture}[scale=1.0]
  % Titulo
  \node[font=\sffamily\bfseries\small, PQCMainBlue] at (5.5, 3.6) {Latencia por operación — ML-KEM-768 vs RSA-2048};
  
  % ML-KEM barra KeyGen
  \fill[PQCTeal!80] (0, 2.5) rectangle (0.8, 3.1);
  \node[right, font=\scriptsize\sffamily] at (0.85, 2.8) {ML-KEM KeyGen: $\approx 0.03$ ms};
  
  % ML-KEM barra Encap
  \fill[PQCTeal!80] (0, 1.8) rectangle (0.4, 2.4);
  \node[right, font=\scriptsize\sffamily] at (0.45, 2.1) {ML-KEM Encapsular: $\approx 0.04$ ms};
  
  % ML-KEM barra Decap
  \fill[PQCTeal!80] (0, 1.1) rectangle (0.5, 1.7);
  \node[right, font=\scriptsize\sffamily] at (0.55, 1.4) {ML-KEM Desencapsular: $\approx 0.04$ ms};
  
  % RSA barra decrypt
  \fill[PQCCrimson!70] (0, 0.2) rectangle (8.5, 0.8);
  \node[right, font=\scriptsize\sffamily] at (8.55, 0.5) {RSA-2048 Descifrar: $\approx 92$ ms};
  
  % Eje x
  \draw[line width=1pt, gray] (0, 0) -- (10, 0);
  \node[font=\scriptsize\sffamily, gray] at (5, -0.3) {Escala relativa de tiempo};
  
\end{tikzpicture}
\caption{ML-KEM es entre 2000 y 3000 veces más rápido que RSA-2048 en las operaciones criptográficas equivalentes.}
\label{fig:rendimiento_kem}
\end{figure}

\begin{ejemplonumerico}[Desencapsulamiento en el proyecto real de votación UMSA]
En el sistema de votación electrónica post-cuántica de la UMSA:
\begin{itemize}
  \item Cada boleta electoral enviada por un ciudadano encapsula una clave con \textbf{ML-KEM-768}.
  \item Al momento de abrir el ánfora electoral, la Autoridad de Conteo ejecuta la operación de desencapsulamiento en apenas \textbf{$0.042$ milisegundos por voto}.
  \item Bajo el estándar RSA-OAEP equivalente, cada voto tardaría $92.137$ ms: exactamente \textbf{$2{,}175\times$ más lento}.
  \item Con $100{,}000$ votos: ML-KEM tarda $\approx 4.2$ segundos; RSA tardaría $\approx 2.5$ horas.
\end{itemize}
\end{ejemplonumerico}

\section{La propiedad de seguridad más importante: CCA2-seguridad}

Los criptógrafos profesionales no se conforman con decir ``parece difícil de romper''. Exigen una garantía matemática formal llamada \textbf{seguridad IND-CCA2} (\textit{Indistinguishability under Chosen Ciphertext Attack 2}).

\begin{intuicion}[La prueba del oráculo maligno]
Imagina que le das al atacante el poder de un oráculo mágico: puede descifrar \textit{cualquier} texto cifrado que quiera, excepto el que precisamente estás tratando de proteger.

Un sistema IND-CCA2-seguro garantiza que incluso con ese poder sobrehumano, el atacante no puede obtener ninguna información útil sobre el mensaje cifrado. ML-KEM satisface esta garantía más fuerte gracias a la transformación \textbf{Fujisaki-Okamoto}, que convierte un esquema base más sencillo en uno completamente robustecido.
\end{intuicion}

Esta propiedad es crítica para aplicaciones reales como el voto electrónico, donde un atacante sofisticado podría intentar enviar mensajes manipulados al sistema para extraer información indirecta sobre las claves.

\begin{resumencapitulo}
\begin{itemize}
  \item ML-KEM encapsula una clave simétrica compacta dentro de una estructura de retículos modulares con ruido, usando $n=256$, $q=3329$, $k=3$.
  \item La clave pública mide apenas 1184 bytes; las operaciones se ejecutan en menos de 0.05 ms.
  \item El encapsulamiento crea un texto cifrado indistinguible de ruido aleatorio sin la clave privada $\mathbf{s}$.
  \item Roberto recupera el secreto mediante cancelación algebraica y redondeo dipolar ($0$ vs. $q/2$).
  \item ML-KEM es formalmente IND-CCA2-seguro bajo la hipótesis de dureza de Module-LWE.
\end{itemize}
\end{resumencapitulo}

\chapter{El Sello Incorruptible: ¿Cómo Funciona ML-DSA?}
\label{cap:mldsa_sello}

\section{El misterio del sello de cera}

En la antigüedad, cuando un rey enviaba un edicto a través de sus mensajeros, no le importaba que el mensaje fuera secreto: quería que todo el pueblo pudiera leerlo en la plaza pública. Lo que realmente necesitaba el rey era una garantía absoluta de que nadie hubiera alterado una sola palabra del texto y que ningún impostor hubiera escrito el mensaje en su nombre.

Para eso utilizaba su anillo y una gota de cera caliente: \textbf{un sello de autenticidad}.

En el mundo digital, la firma digital cumple exactamente ese mismo rol, pero con consecuencias que el rey medieval no podía imaginar: una firma digital matemáticamente correcta es \textbf{verificable por millones de personas simultáneamente}, no puede ser falsificada aunque un atacante tenga una supercomputadora, y no puede ser repudiada por quien la emitió.

\begin{figure}[htbp]
\centering
\begin{tikzpicture}[scale=0.95]
  % Documento
  \node[draw=PQCMainBlue, fill=BoxBgBlue, rounded corners=2mm, line width=1pt, minimum width=3.2cm, minimum height=1.6cm, align=center] (doc) at (0,0) {\textbf{Documento}\\\small (ej. Boleta de voto\\\small o Decreto oficial)};
  
  % Proceso de firma
  \node[draw=PQCPurple, fill=BoxBgPurple, rounded corners=2mm, line width=1pt, minimum width=3.4cm, minimum height=1.6cm, align=center] (sign) at (5,0) {\textbf{Algoritmo ML-DSA}\\\small Usa clave privada $\mathbf{s}$\\\small Genera vector corto $\mathbf{z}$};
  
  % Firma generada
  \node[draw=PQCGreen, fill=BoxBgGreen, rounded corners=2mm, line width=1pt, minimum width=2.8cm, minimum height=1.6cm, align=center] (sig) at (10,0) {\textbf{Firma Digital $\sigma$}\\\small Todos pueden verificar\\\small Nadie puede falsificar};
  
  % Flechas
  \draw[-{Stealth[length=3mm]}, line width=1.5pt, PQCMainBlue] (doc) -- (sign) node[midway, above, font=\footnotesize\sffamily] {Mensaje $\mu$};
  \draw[-{Stealth[length=3mm]}, line width=1.5pt, PQCPurple] (sign) -- (sig) node[midway, above, font=\footnotesize\sffamily] {Sello único};
\end{tikzpicture}
\caption{La firma digital: un sello matemático que vincula de manera inseparable el contenido del documento con la identidad del firmante.}
\label{fig:firma_digital}
\end{figure}

En la nueva era post-cuántica, el estándar mundial oficializado por el NIST bajo la norma \textbf{FIPS 204} es \textbf{ML-DSA} (originalmente conocido como \textit{CRYSTALS-Dilithium}).

\section{¿Qué es una función Hash? La huella digital del documento}

Antes de entender las firmas digitales, necesitamos conocer un ingrediente fundamental: la \textbf{función hash}.

\begin{intuicion}[El compactador matemático de documentos]
Una función hash como SHA3-256 toma un documento de cualquier longitud (¡incluso una novela de 1000 páginas!) y produce un número de exactamente 256 bits, siempre.
\begin{itemize}
  \item Si cambias \textbf{una sola letra} del documento, el hash resultante cambia por completo.
  \item Dado el hash, es imposible reconstruir el documento original.
  \item Dos documentos distintos nunca producen el mismo hash (en la práctica).
\end{itemize}
\end{intuicion}

\begin{figure}[htbp]
\centering
\begin{tikzpicture}[scale=0.92]
  % Documento grande
  \node[draw=PQCMainBlue, fill=BoxBgBlue, rounded corners=2mm, line width=1pt, minimum width=4cm, minimum height=2.2cm, align=center] (doc) at (0,0) {\textbf{Documento}\\\footnotesize ``El votante Juan Pérez\\\footnotesize vota por el candidato A''\\\footnotesize (cualquier longitud)};
  
  % Hash
  \node[draw=PQCAmber, fill=BoxBgAmber, rounded corners=2mm, line width=1pt, minimum width=4.5cm, minimum height=1cm, align=center] (proc) at (6,0) {SHA3-256};
  
  % Resultado (32 bytes fijo)
  \node[draw=PQCCrimson, fill=red!8, rounded corners=2mm, line width=1pt, minimum width=3.8cm, minimum height=1.2cm, align=center, font=\tiny\ttfamily] (h) at (11.5,0) {a3f2...7d9b\\(siempre 256 bits)};
  
  \draw[-{Stealth[length=3mm]}, line width=1.5pt, PQCMainBlue] (doc) -- (proc);
  \draw[-{Stealth[length=3mm]}, line width=1.5pt, PQCAmber] (proc) -- (h) node[midway, above, font=\small\sffamily] {Huella digital};
  
\end{tikzpicture}
\caption{Una función hash comprime cualquier documento a una huella digital de tamaño fijo (256 bits en SHA3-256).}
\label{fig:hash}
\end{figure}

En las firmas digitales, en lugar de firmar el documento completo (que podría ser enorme), siempre se firma el \textbf{hash del documento}. Así, la carga criptográfica es siempre la misma independientemente del tamaño del archivo.

\section{El problema de la aguja en el pajar: Module-SIS}

Para entender cómo se construyen las firmas sobre retículos, debemos conocer al hermano gemelo de LWE: el problema \textbf{SIS} (\textit{Short Integer Solution}, o Solución Entera Corta).

\begin{intuicion}[¿Qué es el problema SIS?]
Imagina que te entregan una matriz pública $\mathbf{A}$ llena de números de $\mathbb{Z}_q$, y te piden resolver la siguiente ecuación homogénea:
\[
\mathbf{A} \cdot \mathbf{z} \equiv \mathbf{0} \pmod q
\]
Si buscas cualquier vector $\mathbf{z}$ cualquiera, cualquier estudiante de primer semestre puede usar eliminación gaussiana y encontrar millones de soluciones con números grandísimos.

Pero ahora viene la regla mágica: \textbf{tienes que encontrar una solución $\mathbf{z}$ donde todas las coordenadas sean números enteros pequeñísimos} (cercanos a cero, sin que $\mathbf{z}$ sea todo ceros).
\end{intuicion}

\begin{figure}[htbp]
\centering
\begin{tikzpicture}[scale=0.9]
  % Ejes
  \draw[->] (-3.5,0) -- (3.5,0) node[right, font=\small] {$z_1$};
  \draw[->] (0,-3.5) -- (0,3.5) node[above, font=\small] {$z_2$};
  
  % Circunferencia de soluciones "pequeñas" (SIS bound)
  \draw[PQCGreen, line width=1.5pt, dashed] (0,0) circle (1.5cm);
  \node[PQCGreen, font=\footnotesize\sffamily] at (2.2, 1.5) {Frontera SIS ($\beta$)};
  
  % Vector solucion corta (buscado)
  \draw[-{Stealth[length=3mm]}, PQCGreen, line width=2pt] (0,0) -- (0.9, 1.1) node[above right, font=\footnotesize\bfseries, PQCGreen] {Solución corta buscada};
  
  % Soluciones largas (fuera del circulo)
  \fill[PQCCrimson!70] (-2.8, 1.5) circle (2.5pt) node[above, font=\tiny, PQCCrimson] {Gauss: $(-2.8, 1.5)$};
  \fill[PQCCrimson!70] (3.0, -2.1) circle (2.5pt);
  \fill[PQCCrimson!70] (-1.9, -2.8) circle (2.5pt);
  \node[PQCCrimson, font=\footnotesize\sffamily] at (2.5, -2.5) {Soluciones largas};
  
\end{tikzpicture}
\caption{El problema SIS: la eliminación gaussiana encuentra fácilmente soluciones grandes (rojo), pero hallar una solución dentro del círculo $\|\mathbf{z}\| \le \beta$ (verde) es computacionalmente infactible.}
\label{fig:sis_geometria}
\end{figure}

Encontrar una combinación no trivial de vectores que sume exactamente cero manteniéndose microscópica es un problema geométrico equivalente a encontrar el vector no nulo más corto de un retículo en cientos de dimensiones (\textit{Shortest Vector Problem}, SVP).

\section{El gran peligro: La fuga del secreto}

Durante años, diseñar firmas con retículos fue un dolor de cabeza para los criptógrafos. ¿Por qué?

Cuando firmas un documento, la idea naive es calcular directamente:
\[
\mathbf{z} = \mathbf{y} + c \cdot \mathbf{s}
\]
donde $\mathbf{y}$ es un vector aleatorio fresco y $c$ es un reto derivado del hash del mensaje.

El problema es que $\mathbf{z}$ depende matemáticamente de $\mathbf{s}$. Aunque individualmente parece inofensivo, la acumulación de muchas firmas revela patrones estadísticos.

\begin{ojo}[¡Cuidado con los promedios estadísticos!]
Si firmas mil documentos distintos, un atacante paciente que recopile todos tus vectores $\mathbf{z}_1, \mathbf{z}_2, \dots, \mathbf{z}_{1000}$ puede calcular el promedio de las formas geométricas. ¡Poco a poco la clave secreta $\mathbf{s}$ empezaría a asomarse como una huella en la arena!

Este ataque se denomina \textbf{ataque por análisis estadístico de firmas} y fue históricamente el principal obstáculo para los esquemas de firmas sobre retículos.
\end{ojo}

\section{La genialidad de Lyubashevsky: El arte de abortar a tiempo}

En 2009, el matemático Vadim Lyubashevsky propuso una solución elegante y definitiva que cambió la historia de la criptografía: la técnica de \textbf{muestreo con rechazo} (\textit{Fiat-Shamir with Aborts}).

\begin{metafora}[El arquero invisible]
Imagina a un arquero secreto que dispara flechas hacia un blanco. Si sopla una ráfaga de viento en una dirección específica, el ángulo del disparo podría revelar desde qué árbol está disparando.

¿Qué hace el arquero inteligente? Cuando siente que el viento es peligroso, \textbf{no dispara}: baja el arco, descarta esa flecha y espera una nueva brisa neutral. Así, cualquier observador que mire las flechas clavadas en el blanco solo verá una distribución geométrica perfectamente simétrica y circular, sin la más mínima pista sobre dónde está oculto el arquero.
\end{metafora}

En ML-DSA, el proceso completo de firma funciona así:

\begin{figure}[htbp]
\centering
\begin{tikzpicture}[scale=0.9]
  % Nodos del algoritmo de firma
  \node[draw=PQCPurple, fill=BoxBgPurple, rounded corners=2mm, align=center, font=\small, minimum width=4cm] (y) at (5, 6) {Genera $\mathbf{y}$ aleatorio\\(vector ``máscara'')};
  \node[draw=PQCMainBlue, fill=BoxBgBlue, rounded corners=2mm, align=center, font=\small, minimum width=4cm] (c) at (5, 4.5) {Calcula $c = \text{Hash}(\mathbf{A}\mathbf{y}, \mu)$\\(el ``reto'' del mensaje)};
  \node[draw=PQCTeal, fill=BoxBgTeal, rounded corners=2mm, align=center, font=\small, minimum width=4cm] (z) at (5, 3) {Calcula $\mathbf{z} = \mathbf{y} + c\mathbf{s}$};
  
  % Decision
  \node[draw=PQCAmber, fill=BoxBgAmber, diamond, aspect=2, font=\small\sffamily, inner sep=2pt] (dec) at (5, 1.5) {¿$\|\mathbf{z}\| \le \beta$?};
  
  % Si ok
  \node[draw=PQCGreen, fill=BoxBgGreen, rounded corners=2mm, align=center, font=\small] (ok) at (9, 1.5) {\textbf{Publicar firma}\\$\sigma = (c, \mathbf{z})$};
  
  % Si falla (abort)
  \node[draw=PQCCrimson, fill=red!8, rounded corners=2mm, align=center, font=\small] (abort) at (1, 1.5) {\textbf{ABORTAR}\\Volver a generar $\mathbf{y}$};
  
  % Flechas
  \draw[-{Stealth[length=2.5mm]}, line width=1.3pt, PQCPurple] (y) -- (c);
  \draw[-{Stealth[length=2.5mm]}, line width=1.3pt, PQCMainBlue] (c) -- (z);
  \draw[-{Stealth[length=2.5mm]}, line width=1.3pt, PQCTeal] (z) -- (dec);
  \draw[-{Stealth[length=2.5mm]}, line width=1.3pt, PQCGreen] (dec) -- node[above, font=\footnotesize] {Sí} (ok);
  \draw[-{Stealth[length=2.5mm]}, line width=1.3pt, PQCCrimson] (dec) -- node[above, font=\footnotesize] {No} (abort);
  \draw[-{Stealth[length=2.5mm]}, line width=1.3pt, PQCCrimson] (abort) -- (1, 6) -- (y);
  
\end{tikzpicture}
\caption{Flujo del algoritmo de firma ML-DSA: cuando la firma candidata es sospechosa, el algoritmo aborta y reintenta hasta que la distribución es estadísticamente pura.}
\label{fig:firma_flujo}
\end{figure}

Concretamente, cuando el firmante calcula el vector candidato $\mathbf{z} = \mathbf{y} + c\mathbf{s}$:
\begin{enumerate}
  \item Mide el tamaño del vector $\mathbf{z}$ (su norma $\|\mathbf{z}\|$).
  \item Si $\mathbf{z}$ cae demasiado cerca de los bordes o contiene cualquier sesgo que dependa de $\mathbf{s}$, el algoritmo simplemente \textbf{aborta el intento}, genera un nuevo $\mathbf{y}$ y vuelve a intentarlo.
  \item Cuando la firma finalmente se publica, su distribución es estadísticamente independiente de la clave secreta. ¡Cero fuga de información!
\end{enumerate}

La clave matemática del truco es el siguiente teorema intuitivo: si elegimos el rango de $\mathbf{y}$ mucho más amplio que el posible efecto de $c\mathbf{s}$, entonces los $\mathbf{z}$ que pasan la prueba tienen exactamente la misma distribución que $\mathbf{y}$ solo. La clave secreta queda estadísticamente invisible.

\section{¿Cómo comprueba cualquier persona la firma?}

La verificación de una firma ML-DSA no requiere conocer ningún secreto. Cualquier persona en el mundo puede comprobar la validez de un voto o un documento siguiendo dos sencillos pasos de álgebra lineal:

\begin{enumerate}
  \item \textbf{Verificación de pequeñez:} Comprueba que las coordenadas del vector firma $\mathbf{z}$ sean efectivamente pequeñas (que no sobrepasen el límite de seguridad $\beta$). Esto garantiza que no fue construida por fuerza bruta.
  \item \textbf{Verificación de coherencia algebraica:} Multiplica la matriz pública por el vector firma y verifica que coincida con el reto $c$ derivado del mensaje:
  \[
  c' = \text{Hash}(\mathbf{Az} - c\mathbf{t}, \mu) \stackrel{?}{=} c
  \]
  donde $\mathbf{t} = \mathbf{As} + \mathbf{e}$ es la clave pública (análoga a $\mathbf{b}$ en ML-KEM).
\end{enumerate}

\begin{intuicion}[¿Por qué funciona la verificación?]
Cuando el firmante calculó $\mathbf{z} = \mathbf{y} + c\mathbf{s}$, entonces:
\[
\mathbf{Az} = \mathbf{A}(\mathbf{y} + c\mathbf{s}) = \mathbf{Ay} + c\mathbf{As}
\]
Y como $\mathbf{t} = \mathbf{As} + \mathbf{e}$, tenemos:
\[
\mathbf{Az} - c\mathbf{t} \approx \mathbf{Ay} - c\mathbf{e}
\]
El verificador solo puede reconstruir el valor $\mathbf{Ay}$ (multiplicado por la máscara aleatoriaoriginal) porque la firma fue construida de manera coherente. Un impostor que no conoce $\mathbf{s}$ no puede construir un $\mathbf{z}$ pequeño que satisfaga simultáneamente estas dos condiciones.
\end{intuicion}

Si ambas condiciones se cumplen, las matemáticas garantizan con rigor absoluto que:
\begin{itemize}
  \item El documento no fue modificado en ni un solo punto (cualquier cambio altera el hash $\mu$).
  \item Quien firmó poseía la clave privada correspondiente (solo él puede construir $\mathbf{z}$).
  \item Ninguna computadora cuántica presente o futura pudo haber falsificado la firma.
\end{itemize}

\section{Tamaños y rendimiento de ML-DSA}

Al igual que ML-KEM, los tamaños reales de ML-DSA resultan sorprendentemente manejables:

\begin{center}
\begin{tabular}{lrrl}
\toprule
\textbf{Componente} & \textbf{ML-DSA-65} & \textbf{ECDSA P-256} & \textbf{Factor} \\
\midrule
Clave pública & 1952 bytes & 64 bytes & $\approx 30\times$ \\
Clave privada & 4032 bytes & 32 bytes & $\approx 126\times$ \\
\textbf{Firma} & \textbf{3309 bytes} & \textbf{72 bytes} & $\approx 46\times$ \\
Tiempo firmar & $\approx 0.18$ ms & $\approx 0.8$ ms & ML-DSA más rápido \\
Tiempo verificar & $\approx 0.14$ ms & $\approx 2.4$ ms & ML-DSA más rápido \\
\bottomrule
\end{tabular}
\end{center}

La firma de ML-DSA es más grande que ECDSA, pero a cambio ofrece \textbf{seguridad cuántica garantizada} y, paradójicamente, es más rápida en la operación de verificación gracias a que la verificación es solo multiplicación matricial (sin costosas operaciones de punto en curva elíptica).

\begin{resumencapitulo}
\begin{itemize}
  \item Las firmas digitales garantizan autenticidad e integridad: que el mensaje proviene del emisor correcto y no fue alterado.
  \item Se firman los \textbf{hashes} de los documentos (huellas digitales compactas de 256 bits), nunca los documentos completos.
  \item ML-DSA se basa en el problema Module-SIS: hallar vectores enteros pequeños que satisfagan ecuaciones lineales modulares.
  \item El \textbf{muestreo con rechazo} de Lyubashevsky elimina cualquier sesgo estadístico que podría filtrar la clave privada.
  \item La verificación es abierta: cualquier ciudadano con la clave pública puede confirmar la autenticidad de un voto o decreto.
\end{itemize}
\end{resumencapitulo}

\chapter{El Candado de Ayer: RSA y por qué las Computadoras Cuánticas lo Rompen}
\label{cap:rsa_y_shor}

\section{El mundo antes de la criptografía de retículos}

Antes de que existieran los retículos como herramienta criptográfica, el mundo digital se protegía con una familia de algoritmos basados en un tipo diferente de problema matemático: la \textbf{factorización de números enteros grandes}.

El representante más famoso de esta familia es \textbf{RSA}, creado en 1977 por Ron Rivest, Adi Shamir y Leonard Adleman (de ahí sus siglas). RSA fue el primer sistema de clave pública ampliamente adoptado y durante décadas protegió las transacciones bancarias, los correos electrónicos y los certificados de identidad digital de todo el mundo.

En este capítulo vamos a aprender cómo funciona RSA con números reales (¡solo necesitas saber dividir!), y luego descubriremos exactamente por qué el algoritmo cuántico de Shor es capaz de destruirlo en cuestión de segundos.

\section{El corazón de RSA: La multiplicación es fácil, la factorización no}

Empieza con un experimento mental:

\begin{center}
\begin{tabular}{ll}
\toprule
\textbf{Tarea} & \textbf{¿Cuánto tardas?} \\
\midrule
Multiplicar $13 \times 17 = ?$ & Menos de 5 segundos \\
Factorizar $221 = ? \times ?$ & Unos minutos (ensayo y error) \\
Multiplicar $4{,}294{,}967{,}311 \times 4{,}294{,}967{,}357 = ?$ & Segundos con calculadora \\
Factorizar $18{,}446{,}744{,}073{,}709{,}551{,}927$ & ¡Horas o días sin ayuda! \\
\bottomrule
\end{tabular}
\end{center}

\begin{metafora}[El producto de primos como una cerradura de combinación]
Imagina que tienes dos números primos $p$ y $q$ (números que solo son divisibles por sí mismos y por 1). Multiplicarlos es tan sencillo como abrir una puerta con su llave: instantáneo. Pero si solo te damos el producto $N = p \times q$ y te pedimos que encuentres de nuevo $p$ y $q$, tienes que probar cientos de miles de posibles divisores uno por uno.

Para números de 2048 bits (que tienen aproximadamente 617 dígitos decimales), el mejor algoritmo clásico de factorización tardaría más tiempo que la edad del universo en encontrar $p$ y $q$.
\end{metafora}

\begin{figure}[htbp]
\centering
\begin{tikzpicture}[scale=0.9]
  % Flecha facil izquierda a derecha
  \node[draw=PQCGreen, fill=BoxBgGreen, rounded corners=2mm, line width=1pt, minimum width=2.5cm, minimum height=1.2cm, align=center] (pq) at (0,0) {$p \times q$\\\small (dos primos)};
  
  \node[draw=PQCMainBlue, fill=BoxBgBlue, rounded corners=2mm, line width=1pt, minimum width=2.5cm, minimum height=1.2cm, align=center] (N) at (5,0) {$N$\\\small (producto público)};
  
  % Facil
  \draw[-{Stealth[length=3mm]}, line width=2pt, PQCGreen] (pq) -- (N) node[midway, above, font=\small\sffamily\bfseries, PQCGreen] {¡Fácil! (milisegundos)};
  
  % Dificil (vuelta)
  \draw[-{Stealth[length=3mm]}, line width=2pt, PQCCrimson, dashed] (N) to[bend left=40] node[midway, below, font=\small\sffamily\bfseries, PQCCrimson, align=center] {¡Difícil para clasicas!\\(siglos de cómputo)} (pq);
  
  % Flecha cuantica
  \draw[-{Stealth[length=3mm]}, line width=2pt, PQCAmber] (N) to[bend right=50] node[midway, above right, font=\small\sffamily\bfseries, PQCAmber, align=center] {Shor cuántico:\\segundos} (pq);
  
\end{tikzpicture}
\caption{La asimetría de la factorización: fácil hacia adelante, imposible hacia atrás para computadoras clásicas, pero trivial para Shor cuántico.}
\label{fig:factorizacion}
\end{figure}

\section{El funcionamiento de RSA paso a paso}

Veamos RSA con números pequeños para visualizar su mecánica real. En la práctica se usan números enormes (2048 bits), pero la estructura matemática es idéntica.

\subsection*{Generación de claves (lo que hace el receptor Roberto)}

\begin{ejemplonumerico}[RSA con números pequeños: $p=11$, $q=13$]
\begin{enumerate}
  \item \textbf{Elige dos primos:} $p = 11$ y $q = 13$.
  
  \item \textbf{Calcula el módulo público:} $N = p \times q = 11 \times 13 = 143$.
  
  \item \textbf{Calcula $\phi(N)$ (función de Euler):}
  \[
  \phi(N) = (p-1)(q-1) = 10 \times 12 = 120
  \]
  Esto cuenta cuántos números entre $1$ y $N$ no comparten factores con $N$.
  
  \item \textbf{Elige el exponente público $e$:} Escoge un número impar pequeño coprimo con $120$. Tomaremos $e = 7$.
  
  \item \textbf{Calcula el exponente privado $d$:} Busca $d$ tal que $e \cdot d \equiv 1 \pmod{\phi(N)}$, es decir, $7d \equiv 1 \pmod{120}$. La solución es $d = 103$ (ya que $7 \times 103 = 721 = 6 \times 120 + 1$).
  
  \item \textbf{Clave Pública:} $(e, N) = (7, 143)$.
  \item \textbf{Clave Privada:} $(d, N) = (103, 143)$.
\end{enumerate}
\end{ejemplonumerico}

\subsection*{Cifrado (lo que hace la emisora Alicia)}

Para cifrar un número $m$ (que representa nuestro mensaje), Alicia usa la clave pública de Roberto:
\[
c = m^e \pmod N = m^7 \pmod{143}
\]

\begin{ejemplonumerico}[Alicia cifra el mensaje $m=9$]
\begin{align*}
c &= 9^7 \pmod{143} \\
9^2 &= 81 \\
9^4 &= 81^2 = 6561 \equiv 6561 - 45 \times 143 = 6561 - 6435 = 126 \pmod{143} \\
9^7 &= 9^4 \cdot 9^2 \cdot 9 = 126 \times 81 \times 9 \pmod{143} \\
126 \times 81 &= 10206 \equiv 10206 - 71 \times 143 = 10206 - 10153 = 53 \pmod{143} \\
53 \times 9 &= 477 \equiv 477 - 3 \times 143 = 477 - 429 = 48 \pmod{143}
\end{align*}
Alicia envía $c = 48$.
\end{ejemplonumerico}

\subsection*{Descifrado (lo que hace Roberto con su clave privada)}

Roberto usa su exponente privado $d$ para recuperar el mensaje:
\[
m = c^d \pmod N = 48^{103} \pmod{143}
\]

Gracias al \textbf{Pequeño Teorema de Fermat} y la estructura del grupo multiplicativo $(\mathbb{Z}_N)^*$, este cálculo siempre devuelve el mensaje original $m = 9$.

\begin{intuicion}[¿Por qué funciona el descifrado?]
El matemático Euler demostró que para cualquier $m$ coprimo con $N = pq$:
\[
m^{\phi(N)} \equiv 1 \pmod N
\]
Por ello, $m^{e \cdot d} = m^{1 + k \cdot \phi(N)} = m \cdot (m^{\phi(N)})^k \equiv m \cdot 1^k = m \pmod N$.

La clave es que calcular $\phi(N) = (p-1)(q-1)$ requiere conocer $p$ y $q$. Un atacante que solo conoce $N$ no puede calcular $\phi(N)$ sin factorizar $N$. Y factorizar $N$ es el problema difícil.
\end{intuicion}

\section{La amenaza cuántica: El algoritmo de Shor}

En 1994, el matemático Peter Shor, trabajando en los Laboratorios Bell, publicó uno de los resultados más impactantes de la historia de la computación: un \textbf{algoritmo cuántico} capaz de factorizar cualquier número entero en tiempo polinomial.

Esto significa que una computadora cuántica suficientemente grande puede romper RSA-2048 en \textbf{horas o incluso minutos}, algo que a una computadora clásica le tomaría millones de años.

\subsection*{¿Qué hace especial a una computadora cuántica?}

\begin{metafora}[El laberinto de infinitos pasillos]
Imagina un laberinto gigante con millones de pasillos ramificados. Para un explorador clásico (una computadora normal), la única manera de encontrar la salida es recorrer los pasillos uno por uno o en pequeños grupos en paralelo.

Una computadora cuántica es como un explorador que puede desdoblarse en \textbf{millones de copias fantasmales simultáneas}, explorando todos los pasillos a la vez. Cuando uno de los fantasmas encuentra la salida, todos los demás colapsan hacia esa solución.
\end{metafora}

El principio físico detrás de esto es la \textbf{superposición cuántica}: a diferencia de los bits clásicos (que son 0 o 1), los \textbf{qubits} cuánticos pueden estar en un estado de superposición donde representan 0 y 1 \textit{al mismo tiempo}, hasta que se miden.

\begin{figure}[htbp]
\centering
\begin{tikzpicture}[scale=0.88]
  % Bit clasico
  \node[draw=gray!60, fill=gray!15, rounded corners=2mm, minimum width=2cm, minimum height=1.2cm, align=center] at (0,0) {\textbf{Bit clásico}\\$0$ \textit{ó} $1$};
  
  % Qubit
  \node[draw=PQCTeal, fill=BoxBgTeal, rounded corners=2mm, minimum width=2.5cm, minimum height=1.6cm, align=center] at (5,0) {\textbf{Qubit cuántico}\\$\alpha|0\rangle + \beta|1\rangle$\\\scriptsize (superposición)};
  
  % n clasicos
  \node[draw=gray!60, fill=gray!15, rounded corners=2mm, minimum width=3cm, minimum height=1.2cm, align=center] at (0,-2.5) {\textbf{$n$ bits clásicos}\\$2^n$ estados posibles,\\solo 1 activo a la vez};
  
  % n qubits
  \node[draw=PQCPurple, fill=BoxBgPurple, rounded corners=2mm, minimum width=3cm, minimum height=1.4cm, align=center] at (5,-2.5) {\textbf{$n$ qubits}\\$2^n$ estados activos\\\textbf{simultáneamente}};
  
  \draw[-{Stealth[length=2.5mm]}, line width=1.3pt, gray!60] (0,0) -- (0,-2.5);
  \draw[-{Stealth[length=2.5mm]}, line width=1.3pt, PQCTeal] (5,0) -- (5,-2.5);
  
\end{tikzpicture}
\caption{La diferencia fundamental: con $n=300$ qubits, una computadora cuántica puede explorar $2^{300}$ estados simultáneamente. Ese número supera la cantidad de átomos en el universo observable.}
\label{fig:clasico_vs_cuantico}
\end{figure}

\subsection*{El corazón de Shor: Detectar periodicidad oculta}

La genialidad del algoritmo de Shor radica en transformar el problema de factorizar $N$ en un problema de encontrar el \textbf{período de una función}, y luego usar el poder cuántico para encontrar ese período instantáneamente.

Aquí está la idea en etapas, sin fórmulas intimidantes:

\begin{figure}[htbp]
\centering
\begin{tikzpicture}[scale=0.9]
  % Columna de pasos
  \node[draw=PQCMainBlue, fill=BoxBgBlue, rounded corners=2mm, line width=1pt, align=center, minimum width=5cm] (s1) at (5, 7) {\textbf{Paso 1: Elegir $a$ al azar}\\\small Escoge un número $a$ con $1 < a < N$};
  
  \node[draw=PQCTeal, fill=BoxBgTeal, rounded corners=2mm, line width=1pt, align=center, minimum width=5cm] (s2) at (5, 5.4) {\textbf{Paso 2: Calcular la función periódica}\\\small $f(x) = a^x \pmod N$};
  
  \node[draw=PQCPurple, fill=BoxBgPurple, rounded corners=2mm, line width=1pt, align=center, minimum width=5.5cm] (s3) at (5, 3.8) {\textbf{Paso 3: Encontrar el período $r$}\\\small tal que $a^r \equiv 1 \pmod N$\\\small {\color{PQCGreen}$\Leftarrow$ ¡Aquí usa el poder cuántico!}};
  
  \node[draw=PQCGreen, fill=BoxBgGreen, rounded corners=2mm, line width=1pt, align=center, minimum width=5cm] (s4) at (5, 2.1) {\textbf{Paso 4: Calcular el factor}\\\small $\gcd(a^{r/2} - 1, N)$ = factor de $N$};
  
  \node[draw=PQCCrimson, fill=red!8, rounded corners=2mm, line width=1pt, align=center, minimum width=4cm] (s5) at (5, 0.5) {\textbf{¡RSA destruido!}\\\small Se conocen $p$ y $q$};
  
  \draw[-{Stealth[length=2.5mm]}, line width=1.3pt] (s1) -- (s2);
  \draw[-{Stealth[length=2.5mm]}, line width=1.3pt] (s2) -- (s3);
  \draw[-{Stealth[length=2.5mm]}, line width=1.3pt] (s3) -- (s4);
  \draw[-{Stealth[length=2.5mm]}, line width=1.3pt] (s4) -- (s5);
  
\end{tikzpicture}
\caption{Flujo simplificado del algoritmo de Shor. El paso crítico (detectar el período $r$) es exponencialmente difícil para computadoras clásicas pero se resuelve en tiempo polinomial con qubits cuánticos.}
\label{fig:shor_flujo}
\end{figure}

\subsection*{¿Por qué encontrar el período destruye RSA?}

Cuando encuentras el período $r$ de la función $a^x \pmod N$, ocurre algo algebraicamente hermoso. Si $r$ es par:
\[
a^r \equiv 1 \pmod N \quad \Rightarrow \quad (a^{r/2} - 1)(a^{r/2} + 1) \equiv 0 \pmod N
\]

Esto significa que $N$ divide el producto, pero no necesariamente cada factor individual. Al calcular el máximo común divisor:
\[
\gcd(a^{r/2} - 1, N) \quad \text{y} \quad \gcd(a^{r/2} + 1, N)
\]
con altísima probabilidad obtienes directamente $p$ y $q$. ¡El candado RSA queda abierto!

\begin{ejemplonumerico}[Shor con $N = 143$, $a = 2$]
Calculamos la secuencia $2^x \pmod{143}$:
\begin{center}
\begin{tabular}{cccccccc}
$x$: & 1 & 2 & 3 & 4 & $\cdots$ & 59 & 60 \\
$2^x$: & 2 & 4 & 8 & 16 & $\cdots$ & 33 & 66 \\
$x$: & 61 & 62 & $\cdots$ & \textbf{120} & & & \\
$2^x$: & 132 & 121 & $\cdots$ & \textbf{1} & & & \\
\end{tabular}
\end{center}
El período es $r = 120$. Como $r$ es par, calculamos:
\[
\gcd(2^{60} - 1, 143) = \gcd(\ldots, 143)
\]
Aplicando el algoritmo de Euclides se obtiene $13$. Y en efecto: $143 = 11 \times 13$.
\end{ejemplonumerico}

\section{¿Por qué una computadora cuántica encuentra el período tan rápido?}

Una computadora clásica necesitaría calcular $2^1, 2^2, 2^3, \dots$ hasta encontrar que $2^{120} \equiv 1$: hasta $120$ evaluaciones para este ejemplo pequeño, pero para un $N$ de 2048 bits serían potencialmente $2^{1024}$ evaluaciones.

La computadora cuántica aplica la \textbf{Transformada Cuántica de Fourier} (QFT), un equivalente cuántico de la Transformada Rápida de Fourier (FFT) que ya conoces de señales y audio:

\begin{figure}[htbp]
\centering
\begin{tikzpicture}[scale=0.92]
  % Funcion periodica (onda)
  \begin{scope}[shift={(0,0)}]
    \draw[line width=1.5pt, PQCTeal, domain=0:8, samples=200] plot (\x, {0.6*sin(\x*180/3.14*2/4)+1.5});
    \node[font=\small\sffamily, PQCTeal] at (4, 2.4) {$f(x) = a^x \pmod N$ (periódica con período $r$)};
    \draw[dashed, PQCAmber, line width=1pt] (0,0.3) -- (8,0.3);
    \draw[{Stealth[length=2mm]}-{Stealth[length=2mm]}, PQCAmber, line width=1pt] (0,0.3) -- (4,0.3) node[below, font=\tiny] {período $r$};
  \end{scope}
  
  % QFT
  \node[draw=PQCPurple, fill=BoxBgPurple, rounded corners=3mm, minimum width=2cm, minimum height=1cm, font=\small\sffamily\bfseries] (qft) at (9, 1.5) {QFT cuántica};
  
  % Espectro
  \begin{scope}[shift={(11,0)}]
    \foreach \x/\h in {1/0.2, 2/0.2, 3/2.2, 4/0.2, 5/0.2, 6/2.1, 7/0.2} {
      \fill[PQCGreen!80] ({\x*0.9-0.45}, 0) rectangle ({\x*0.9}, \h);
    }
    \node[font=\small\sffamily, PQCGreen] at (3.15, 2.7) {Picos = $N/r$};
  \end{scope}
  
  % Flecha QFT
  \draw[-{Stealth[length=2.5mm]}, line width=1.5pt, PQCPurple] (8.2, 1.5) -- (qft);
  \draw[-{Stealth[length=2.5mm]}, line width=1.5pt, PQCPurple] (qft) -- (10.1, 1.5);
  
\end{tikzpicture}
\caption{La QFT cuántica detecta el período $r$ de $f(x)$ en un solo paso: busca los picos de frecuencia en el espectro, que están separados exactamente por $N/r$.}
\label{fig:qft}
\end{figure}

La QFT coloca todos los $2^n$ posibles valores de $x$ en superposición simultánea, evalúa $f(x)$ en todos ellos a la vez, y luego extrae la información de periodicidad mediante interferencia cuántica. El período $r$ emerge de los picos del espectro.

Para un número de 2048 bits, la QFT necesita apenas del orden de $(2048)^2 \approx 4 \times 10^6$ operaciones cuánticas, en lugar de los $2^{1024}$ pasos de un algoritmo clásico.

\section{¿Y los retículos? La inmunidad a Shor}

La pregunta natural es: ¿por qué el algoritmo de Shor no rompe también a ML-KEM y ML-DSA?

\begin{intuicion}[La diferencia fundamental entre RSA y retículos]
\begin{itemize}
  \item \textbf{RSA} esconde el secreto en la \textbf{estructura periódica} del grupo multiplicativo $(\mathbb{Z}_N)^*$. La función $a^x \pmod N$ es perfectamente periódica, limpia, como una onda de radio pura. La QFT detecta esa onda en un instante.
  
  \item \textbf{Module-LWE} esconde el secreto en \textbf{datos ruidosos}: $\mathbf{b} = \mathbf{As} + \mathbf{e}$. El error $\mathbf{e}$ destruye completamente cualquier estructura periódica. No hay onda, solo caos aleatorio. La QFT no tiene nada que sintonizar.
\end{itemize}
\end{intuicion}

\begin{figure}[htbp]
\centering
\begin{tikzpicture}[scale=0.92]
  % RSA: onda limpia
  \begin{scope}[shift={(0,2.5)}]
    \node[font=\small\sffamily\bfseries, PQCCrimson] at (4, 2) {\textbf{RSA}: función periódica limpia $\Rightarrow$ Shor la detecta};
    \draw[line width=1.8pt, PQCCrimson, domain=0:8, samples=200] plot (\x, {0.7*sin(\x*180/3.14*2/2)});
    \draw[dashed, PQCAmber, line width=1pt] (0,-0.9) -- (8,-0.9);
    \draw[{Stealth[length=2mm]}-{Stealth[length=2mm]}, PQCAmber] (0,-0.9) -- (2,-0.9) node[midway, below, font=\tiny] {período};
  \end{scope}
  
  % Ruido LWE
  \begin{scope}[shift={(0,0)}]
    \node[font=\small\sffamily\bfseries, PQCGreen] at (4, 1.2) {\textbf{LWE}: datos ruidosos sin período $\Rightarrow$ Shor no tiene nada};
    \foreach \x in {0.3,0.6,...,7.8} {
      \pgfmathsetmacro\yr{(rand)*0.7}
      \draw[line width=1.5pt, PQCGreen!70] (\x, 0) -- (\x, \yr);
    }
  \end{scope}
  
\end{tikzpicture}
\caption{La razón geométrica de la inmunidad post-cuántica: RSA tiene una señal periódica limpia que Shor detecta; LWE tiene ruido puro sin patrón que ningún algoritmo cuántico puede explotar.}
\label{fig:rsa_vs_lwe_shor}
\end{figure}

Formalmente, los mejores algoritmos cuánticos conocidos para atacar problemas de retículos (como BKZ cuántico) siguen siendo \textbf{exponenciales} en la dimensión del retículo. Para dimensiones de $n = 768$ como las usadas en ML-KEM-768, la complejidad de romper el sistema se estima en $2^{178}$ operaciones cuánticas, lo cual está firmemente fuera del alcance de cualquier tecnología imaginable.

\section{El estado actual de las computadoras cuánticas}

En 2024, las computadoras cuánticas más avanzadas tienen entre 400 y 1000 qubits \textbf{físicos}, pero con tasas de error tan altas que en la práctica equivalen a solo unos pocos \textbf{qubits lógicos} estables. Para ejecutar el algoritmo de Shor sobre RSA-2048 se estima que se necesitan aproximadamente \textbf{4 millones de qubits lógicos} sin errores.

\begin{intuicion}[¿Por qué migrar ahora si los qubits no existen aún?]
El concepto de \textbf{``cosechar ahora, descifrar después''} (\textit{Harvest now, decrypt later}) es una amenaza real: adversarios sofisticados pueden estar capturando y almacenando hoy el tráfico cifrado con RSA, esperando el día en que tengan suficientes qubits para descifrarlo retroactivamente.

Datos cifrados hoy con RSA-2048 que deban mantenerse secretos por 20 o 30 años (como archivos médicos, estrategias de estado o claves de infraestructura crítica) están en riesgo si no se migran ahora a algoritmos post-cuánticos como ML-KEM y ML-DSA.
\end{intuicion}

Por eso el NIST publicó sus estándares post-cuánticos en 2024, y por eso el proyecto de grado de la UMSA implementa ML-KEM y ML-DSA en lugar de RSA para el sistema de votación electrónica.

\begin{resumencapitulo}
\begin{itemize}
  \item RSA se basa en la dificultad de \textbf{factorizar} el producto de dos números primos grandes.
  \item El proceso de cifrado/descifrado RSA usa potenciación modular: $c = m^e \pmod N$ y $m = c^d \pmod N$.
  \item El algoritmo de Shor transforma la factorización en detección de periodicidad, usando la QFT cuántica para encontrar el período $r$ de $a^x \pmod N$ en tiempo polinomial.
  \item LWE/Module-LWE no tiene estructura periódica (el ruido $\mathbf{e}$ la destruye), por lo que Shor no puede aplicarse.
  \item La migración a criptografía post-cuántica es urgente hoy debido al riesgo de ``cosechar ahora, descifrar después''.
\end{itemize}
\end{resumencapitulo}

\chapter{Epílogo y Glosario Amigable}
\label{cap:epilogo_glosario}

\section{Epílogo: La matemática viva que protege al mundo}

¡Felicidades por llegar hasta aquí!

Al iniciar este libro, tal vez pensabas que la criptografía post-cuántica era un territorio impenetrable, reservado únicamente para matemáticos con décadas de experiencia o investigadores de laboratorios nucleares.

Sin embargo, a lo largo de estas páginas has descubierto una verdad inspiradora:
\begin{enumerate}
  \item Los conceptos fundamentales no son monstruos incomprensibles; son \textbf{ideas geométricas puras y elegantes}.
  \item Una cuadrícula de puntos regulares en el plano que aprendiste a graficar en tu cuaderno de colegio es el germen de un \textbf{retículo multidimensional}.
  \item La aritmética de las horas de un reloj que utilizas para saber cuándo empieza tu clase es la misma \textbf{aritmética modular $\mathbb{Z}_q$} que confina los secretos digitales.
  \item El juego de mezclar pintura o susurrar en un concierto ruidoso es exactamente lo que hace el problema \textbf{LWE}, donde un error deliberado frustra los métodos convencionales de resolución.
  \item Los polinomios que sumas en el pizarrón permiten que las matrices viajen livianas por las redes de fibra óptica en \textbf{Module-LWE}.
\end{enumerate}

Las dos grandes primitivas que acabas de comprender ---\textbf{ML-KEM} para el candado de confidencialidad y \textbf{ML-DSA} para el sello de autenticidad--- no son teorías de ciencia ficción archivadas en una biblioteca: son los estándares oficiales del Instituto Nacional de Estándares y Tecnología (NIST) aprobados en agosto de 2024 para blindar toda la infraestructura informática del planeta durante los próximos cincuenta años.

Y son exactamente las primitivas matemáticas que dan vida al artefacto de investigación desarrollado en la Universidad Mayor de San Andrés (UMSA) para garantizar que los votos ciudadanos sean anónimos, incorruptibles y eternamente seguros frente a cualquier máquina cuántica del futuro.

\section{Glosario Amigable para Primer Semestre}

Guarda este glosario como una brújula rápida para tus estudios:

\begin{description}
  \item[\textbf{Aritmética Modular ($\mathbb{Z}_q$):}] Sistema numérico circular finito donde los cálculos dan la vuelta al llegar a un valor tope $q$, quedándose con el residuo de la división entera.
  
  \item[\textbf{Base de un Retículo:}] Conjunto mínimo de vectores cuyas combinaciones lineales enteras generan todos y cada uno de los puntos de la red. Una base ``buena'' tiene vectores cortos y casi ortogonales; una base ``mala'' tiene vectores largos y distorsionados.
  
  \item[\textbf{CVP (\textit{Closest Vector Problem}):}] Problema geométrico que consiste en encontrar el punto del retículo más cercano a un punto que fue perturbado por un pequeño ruido en el espacio.
  
  \item[\textbf{Firma Digital:}] Equivalente criptográfico moderno de un sello de cera real: garantiza que el mensaje no fue modificado (integridad) y que proviene legítimamente del autor (autenticidad).
  
  \item[\textbf{KEM (\textit{Key Encapsulation Mechanism}):}] Protocolo criptográfico diseñado para acordar una clave secreta simétrica compacta entre dos personas a través de un canal público no confiable.
  
  \item[\textbf{LWE (\textit{Learning With Errors}):}] Problema computacional introducido por Oded Regev en 2005. Consiste en resolver un sistema lineal ruidoso $\mathbf{b} = \mathbf{A}\mathbf{s} + \mathbf{e} \pmod q$. Es inmune a computadoras clásicas y cuánticas.
  
  \item[\textbf{ML-DSA (FIPS 204):}] Algoritmo de firma digital post-cuántico basado en retículos modulares (Module-LWE y Module-SIS). Anteriormente conocido como *CRYSTALS-Dilithium*.
  
  \item[\textbf{ML-KEM (FIPS 203):}] Mecanismo de encapsulamiento de claves post-cuántico basado en Module-LWE. Anteriormente conocido como *CRYSTALS-Kyber*.
  
  \item[\textbf{Module-LWE (M-LWE):}] Variante algebraica de LWE donde las coordenadas de los vectores y las entradas de las matrices son polinomios de un anillo $R_q = \mathbb{Z}_q[x]/(x^n+1)$, ofreciendo tamaño compacto y cálculo ultrasónico.
  
  \item[\textbf{Muestreo con Rechazo (\textit{Rejection Sampling}):}] Técnica estadística empleada en firmas digitales de retículos donde los intentos que podrían filtrar información sobre la clave secreta se descartan y recalculan antes de publicarse.
  
  \item[\textbf{Retículo (\textit{Lattice}):}] Conjunto discreto e infinito de puntos en un espacio vectorial generado por todas las combinaciones lineales enteras de un conjunto de vectores base linealmente independientes.
  
  \item[\textbf{SIS (\textit{Short Integer Solution}):}] Problema de hallar un vector no nulo de números enteros muy pequeños $\mathbf{z}$ tal que $\mathbf{A}\mathbf{z} \equiv \mathbf{0} \pmod q$.
\end{description}

\vspace{0.8cm}
\begin{tcolorbox}[colback=BoxBgBlue, colframe=PQCMainBlue, arc=2mm, title={\sffamily\bfseries Para Seguir Aprendiendo}]
\sffamily\small
Si quieres profundizar en este fascinante viaje:
\begin{itemize}
  \item Revisa las especificaciones oficiales del NIST: \textbf{FIPS 203} (ML-KEM) y \textbf{FIPS 204} (ML-DSA) publicadas en 2024.
  \item Consulta los artículos fundacionales de Oded Regev (LWE, 2005) y Vadim Lyubashevsky (Firmas con rechazo, 2009).
  \item Explora el código fuente y benchmarks del proyecto de grado en el repositorio \texttt{src/pqc} de la Carrera de Informática de la UMSA.
\end{itemize}
\end{tcolorbox}


\end{document}
